\begin{proof}[Proof of \cref{obj:prop:signed-causal-bridge}]
Throughout, \(d\ge2\), so every stratum has positive probability under a law in \(\mathcal V_d\).

\begin{enumerate}
\item \emph{Identification of the signed mean.}
Fix \(P\in\mathcal V_d\) and \(j\in[d]\). By \cref{obj:ass:fair-randomization}, summing over the potential outcome for the other arm gives, for \(a\in\{0,1\}\),
\[
\begin{aligned}
&P(X=j,A=a,Y^a=1)\\
&\qquad=\sum_{y\in\{0,1\}}
 P(X=j,A=a,Y^a=1,Y^{1-a}=y)\\
&\qquad=\frac12\sum_{y\in\{0,1\}}
 P(X=j,Y^a=1,Y^{1-a}=y)
 =\frac12P(X=j,Y^a=1).
\end{aligned}
\]
Summing the analogous identity over both values of \(Y^a\) yields
\[
P(X=j,A=a)=\frac12P(X=j)=\frac1{2d},
\]
where the last equality uses \cref{obj:ass:uniform-covariates}.

By \cref{obj:ass:consistency}, the following indicator identity holds almost surely:
\[
\begin{aligned}
\mathbf1_{\{X=j\}}S
={}&2\mathbf1_{\{X=j,A=1,Y^1=1\}}
-\mathbf1_{\{X=j,A=1\}}\\
&-2\mathbf1_{\{X=j,A=0,Y^0=1\}}
+\mathbf1_{\{X=j,A=0\}}.
\end{aligned}
\]
Integrating and using the fair-assignment identities gives
\[
\mathbb E_P[\mathbf1_{\{X=j\}}S]
=P(X=j,Y^1=1)-P(X=j,Y^0=1).
\]
Dividing by \(P(X=j)=1/d>0\), we obtain
\[
\mathbb E_P[S\mid X=j]
=\mu_{1j}(P)-\mu_{0j}(P)
=\tau_j(P).
\]
% lean: CausalSmith.Stat.LdpOptvalueUniformFrontier.signedCellMean_eq_contrast

\item \emph{Identification of the paired law.}
Because \(S\) takes values in \(\{-1,1\}\), its two signed atoms satisfy
\[
\begin{aligned}
P(X=j,S=1)+P(X=j,S=-1)&=\frac1d,\\
P(X=j,S=1)-P(X=j,S=-1)
&=\mathbb E_P[\mathbf1_{\{X=j\}}S]
=\frac{\tau_j(P)}d.
\end{aligned}
\]
Solving these two equations gives, for \(s\in\{-1,1\}\),
\[
P(X=j,S=s)=\frac{1+s\tau_j(P)}{2d}.
\]
Moreover, \cref{obj:ass:interior-means} gives
\[
-\frac12
=\frac14-\frac34
\le \tau_j(P)
\le \frac34-\frac14
=\frac12.
\]
For any \(\theta\in[-1/2,1/2]^d\), the paired masses satisfy
\[
\frac{1+s\theta_j}{2d}\ge0,
\qquad
\sum_{j=1}^d\sum_{s\in\{-1,1\}}
\frac{1+s\theta_j}{2d}
=\sum_{j=1}^d\frac1d=1.
\]
Thus \(P_{\tau(P)}\) is a probability law. Equality of the displayed atom probabilities on the finite alphabet proves
\[
\mathcal L_P(X,S)=P_{\tau(P)}.
\]
% lean: CausalSmith.Stat.LdpOptvalueUniformFrontier.causal_signed_law

\item \emph{Decomposition of the causal value.}
For each stratum, the two cases
\(\mu_{0j}(P)\le\mu_{1j}(P)\) and
\(\mu_{1j}(P)\le\mu_{0j}(P)\) give
\[
\max\{\mu_{0j}(P),\mu_{1j}(P)\}
=\frac{\mu_{0j}(P)+\mu_{1j}(P)}2
+\frac{|\tau_j(P)|}2.
\]
Consequently,
\[
V(P)
=\frac1d\sum_{j=1}^d
 \frac{\mu_{0j}(P)+\mu_{1j}(P)}2
+\frac{F(\tau(P))}2.
\]

To identify the first term, consistency gives the almost-sure identity
\[
\mathbf1_{\{X=j\}}Y
=\mathbf1_{\{X=j,A=0,Y^0=1\}}
+\mathbf1_{\{X=j,A=1,Y^1=1\}}.
\]
The fair-assignment identity
\(P(X=j,A=a,Y^a=1)=\frac12P(X=j,Y^a=1)\)
therefore yields
\[
\begin{aligned}
\mathbb E_P[\mathbf1_{\{X=j\}}Y]
&=\frac{P(X=j,Y^0=1)+P(X=j,Y^1=1)}2\\
&=\frac1d\,\frac{\mu_{0j}(P)+\mu_{1j}(P)}2.
\end{aligned}
\]
Since the strata partition the sample space,
\[
B(P)=\mathbb E_P[Y]
=\sum_{j=1}^d\mathbb E_P[\mathbf1_{\{X=j\}}Y]
=\frac1d\sum_{j=1}^d
 \frac{\mu_{0j}(P)+\mu_{1j}(P)}2.
\]
Substitution proves
\[
V(P)=B(P)+\frac{F(\tau(P))}2.
\]
% lean: CausalSmith.Stat.LdpOptvalueUniformFrontier.causal_value_decomposition

\item \emph{Total variation from the uniform law.}
Fix \(\theta\in[-1/2,1/2]^d\). Both \(P_\theta\) and
\(\mathsf U_{2d}=P_0\) are probability laws by the mass calculation above.
Define the event of nonnegative atom differences by
\[
\mathcal E_{\theta,+}
:=\left\{(j,s)\in[d]\times\{-1,1\}:
P_\theta(j,s)-\frac1{2d}\ge0\right\}.
\]
The atom differences sum to zero:
\[
\sum_{j=1}^d\sum_{s\in\{-1,1\}}
\left(P_\theta(j,s)-\frac1{2d}\right)=0.
\]
Splitting this sum into its nonnegative and negative terms gives
\[
P_\theta(\mathcal E_{\theta,+})
-\mathsf U_{2d}(\mathcal E_{\theta,+})
=\frac12\sum_{j=1}^d\sum_{s\in\{-1,1\}}
\left|P_\theta(j,s)-\frac1{2d}\right|.
\]
For every event \(\mathcal E\subseteq[d]\times\{-1,1\}\), its sum of atom
differences lies between the sum of all negative differences and the sum
of all nonnegative differences. Hence
\[
\left|P_\theta(\mathcal E)-\mathsf U_{2d}(\mathcal E)\right|
\le\frac12\sum_{j=1}^d\sum_{s\in\{-1,1\}}
\left|P_\theta(j,s)-\frac1{2d}\right|.
\]
The event \(\mathcal E_{\theta,+}\) attains this bound, proving the finite
atom formula for total variation. For each \(j\), its two summands reduce to
\[
\begin{aligned}
\sum_{s\in\{-1,1\}}
\left|\frac{1+s\theta_j}{2d}-\frac1{2d}\right|
&=\left|\frac{\theta_j}{2d}\right|
+\left|-\frac{\theta_j}{2d}\right|\\
&=\frac{|\theta_j|}{d}.
\end{aligned}
\]
Therefore
\[
\operatorname{TV}(P_\theta,\mathsf U_{2d})
=\frac1{2d}\sum_{j=1}^d|\theta_j|
=\frac{F(\theta)}2.
\]
Applying this identity at \(\theta=\tau(P)\) proves
\[
V(P)=B(P)+\operatorname{TV}(P_{\tau(P)},\mathsf U_{2d}).
\]
% lean: CausalSmith.Stat.LdpOptvalueUniformFrontier.tvFromUniform_pairedLaw

\item \emph{The symmetric causal construction.}
Fix \(\theta\in[-1/2,1/2]^d\). For
\(a,y\in\{0,1\}\) and \(j\in[d]\), define the Bernoulli masses
\[
\pi_{aj}(y)
:=
\left(\frac{1+(2a-1)\theta_j}{2}\right)^y
\left(\frac{1-(2a-1)\theta_j}{2}\right)^{1-y}.
\]
Their success probabilities belong to \([1/4,3/4]\), so
\[
\pi_{aj}(y)\ge0,
\qquad
\sum_{y\in\{0,1\}}\pi_{aj}(y)=1.
\]
By \cref{obj:def:symmetric-law}, the full-data atom masses are
\[
\begin{aligned}
&G_\theta(X=j,A=a,Y=y,Y^0=y^0,Y^1=y^1)\\
&\qquad=\frac1{2d}\,
\pi_{0j}(y^0)\pi_{1j}(y^1)
\mathbf1_{\{y=(1-a)y^0+ay^1\}},
\qquad a,y,y^0,y^1\in\{0,1\}.
\end{aligned}
\]
Summing first over \(y\), then over the two potential outcomes and
assignment, gives total mass
\[
\sum_{j=1}^d\sum_{a\in\{0,1\}}\frac1{2d}
\left(\sum_{y^0\in\{0,1\}}\pi_{0j}(y^0)\right)
\left(\sum_{y^1\in\{0,1\}}\pi_{1j}(y^1)\right)
=1.
\]
Thus \(G_\theta\) is a probability law. The same atom formula gives
\[
G_\theta(X=j)=\frac1d
\]
and
\[
\begin{aligned}
&G_\theta(X=j,A=a,Y^0=y^0,Y^1=y^1)\\
&\qquad=\frac1{2d}\pi_{0j}(y^0)\pi_{1j}(y^1)
=\frac12G_\theta(X=j,Y^0=y^0,Y^1=y^1).
\end{aligned}
\]
These are the conditions in
\cref{obj:ass:uniform-covariates,obj:ass:fair-randomization}.
The atom formula assigns all mass to \(Y=Y_A^{\mathrm{pot}}\), supplying
\cref{obj:ass:consistency}. Finally, summing over the other coordinates yields
\[
G_\theta(X=j,Y^a=1)
=\frac1d\,\frac{1+(2a-1)\theta_j}{2},
\]
and hence
\[
\mu_{0j}(G_\theta)=\frac{1-\theta_j}{2},
\qquad
\mu_{1j}(G_\theta)=\frac{1+\theta_j}{2}.
\]
Both means lie in \([1/4,3/4]\), as required by
\cref{obj:ass:interior-means}. All four conditions of
\cref{obj:def:causal-model} therefore hold.
% lean: CausalSmith.Stat.LdpOptvalueUniformFrontier.symmetricLaw_causalModel

\item \emph{Symmetric value and the ancillary assignment.}
The arm means just computed satisfy
\[
\tau_j(G_\theta)=\theta_j,
\qquad
\frac{\mu_{0j}(G_\theta)+\mu_{1j}(G_\theta)}2=\frac12.
\]
Applying the armwise maximum decomposition gives
\[
V(G_\theta)
=\frac1d\sum_{j=1}^d\frac12+\frac{F(\theta)}2
=\frac12+\frac{F(\theta)}2.
\]
% lean: CausalSmith.Stat.LdpOptvalueUniformFrontier.symmetricLaw_value

Define the observed-record reconstruction map on the whole paired
alphabet and both assignment values by
\[
\operatorname{rec}(j,s,a)
:=\left(j,a,\frac{1+s(2a-1)}2\right),
\qquad
j\in[d],\quad s\in\{-1,1\},\quad a\in\{0,1\}.
\]
Its outcome coordinate is binary. Summing the full-data atom masses over
the unused potential outcome gives
\[
\begin{aligned}
G_\theta(X=j,S=s,A=a)
&=\frac1{2d}\,
\pi_{aj}\!\left(\frac{1+s(2a-1)}2\right)\\
&=\frac1{2d}\,\frac{1+s\theta_j}{2}
=\frac{P_\theta(j,s)}2.
\end{aligned}
\]
Indeed, for \(a=0\), the reconstructed outcome is \((1-s)/2\), so the
two choices of \(s\) select the failure and success masses of
\(\operatorname{Bernoulli}((1-\theta_j)/2)\). For \(a=1\), it is
\((1+s)/2\), selecting the success and failure masses of
\(\operatorname{Bernoulli}((1+\theta_j)/2)\). In both cases the selected
mass is \((1+s\theta_j)/2\). Thus the joint signed-input and assignment
law is
\[
\mathcal L_{G_\theta}(X,S,A)
=P_\theta\otimes\operatorname{Bernoulli}(1/2).
\]
% lean: CausalSmith.Stat.LdpOptvalueUniformFrontier.symmetricLaw_ancillary

For every binary observed record, \((2A-1)^2=1\), and therefore
\[
S(2A-1)
=(2A-1)(2Y-1)(2A-1)=2Y-1.
\]
This identity holds pointwise and hence \(G_\theta\)-almost surely.
% lean: CausalSmith.Stat.LdpOptvalueUniformFrontier.outcome_sign_reconstruction

\item \emph{Class preservation by averaging.}
Fix \(\varepsilon\in\mathbb R\) and an original-record protocol \(Q\).
Using the same seed law and message spaces, define
\[
\mathsf K_i(\cdot\mid(j,s),h,r)
:=\frac12\sum_{a\in\{0,1\}}
Q_i(\cdot\mid\operatorname{rec}(j,s,a),h,r).
\]
Each kernel is a finite mixture of measurable probability kernels, so
this defines a protocol.

Suppose \(Q\) satisfies \cref{obj:ass:full-record-privacy}. For every
\(i,h,r\), measurable message event \(E\), and paired inputs
\((j,s),(j',s')\), apply that inequality separately with assignment zero
and assignment one:
\[
Q_i(E\mid\operatorname{rec}(j,s,a),h,r)
\le e^\varepsilon
Q_i(E\mid\operatorname{rec}(j',s',a),h,r),
\qquad a\in\{0,1\}.
\]
Adding the two inequalities with weights \(1/2\) gives
\[
\mathsf K_i(E\mid(j,s),h,r)
\le e^\varepsilon\mathsf K_i(E\mid(j',s'),h,r).
\]
This proves sequential membership as specified by
\cref{obj:def:sequential-class}. If \(Q\) also satisfies
\cref{obj:ass:noninteraction}, then for all histories \(h,h'\),
\[
\begin{aligned}
\mathsf K_i(\cdot\mid(j,s),h,r)
&=\frac12\sum_a
Q_i(\cdot\mid\operatorname{rec}(j,s,a),h,r)\\
&=\frac12\sum_a
Q_i(\cdot\mid\operatorname{rec}(j,s,a),h',r)
=\mathsf K_i(\cdot\mid(j,s),h',r).
\end{aligned}
\]
Thus \cref{obj:def:noninteractive-class} is preserved as well. These
arguments apply at every declared history and seed, for every real
\(\varepsilon\).
% lean: CausalSmith.Stat.LdpOptvalueUniformFrontier.signed_protocol_classes

\item \emph{Reconstruction of the complete sampled input.}
For deterministic vectors, define
\[
\begin{gathered}
\mathbf v=((j_i,s_i))_{i=1}^n
\in([d]\times\{-1,1\})^n,
\qquad
\mathbf a=(a_i)_{i=1}^n\in\{0,1\}^n,\\
\operatorname{rec}^{(n)}(\mathbf v,\mathbf a)
:=\bigl(\operatorname{rec}(j_i,s_i,a_i)\bigr)_{i=1}^n.
\end{gathered}
\]
The single-participant ancillary identity and
\cref{obj:ass:iid-people} imply
\[
\mathcal L_{G_\theta}\!\left(
((X_i,S_i))_{i=1}^n,(A_i)_{i=1}^n\right)
=P_\theta^{\otimes n}
 \otimes\operatorname{Bernoulli}(1/2)^{\otimes n}.
\]
To see the factorization directly, the probability of a specified pair
of vectors is
\[
\prod_{i=1}^n\frac{P_\theta(j_i,s_i)}2
=\left(\prod_{i=1}^nP_\theta(j_i,s_i)\right)2^{-n}.
\]
Reconstruction recovers every observed record, so pushing this product
law through \(\operatorname{rec}^{(n)}\) gives
\[
\bigl(P_\theta^{\otimes n}
 \otimes\operatorname{Bernoulli}(1/2)^{\otimes n}\bigr)
 \circ(\operatorname{rec}^{(n)})^{-1}
=(G_\theta)_O^{\otimes n}.
\]

Write \(\lambda=\mathcal L(R)\) for the protocol's seed law. Including
the independent randomness from
\cref{obj:ass:independent-randomness}, define
\[
\Lambda_\theta
:=P_\theta^{\otimes n}
 \otimes\operatorname{Bernoulli}(1/2)^{\otimes n}
 \otimes\lambda\otimes\operatorname{Uniform}[0,1]
\]
on tuples \((\mathbf v,\mathbf a,r,u)\). Its pushforward satisfies
\[
\mathcal L_{G_\theta}(\mathbf O,R,U)
=\Lambda_\theta\circ
\bigl[(\mathbf v,\mathbf a,r,u)
\mapsto(\operatorname{rec}^{(n)}(\mathbf v,\mathbf a),r,u)\bigr]^{-1}.
\]
Thus running \(Q\) on the sampled causal records has the same joint
message and randomizer law as running it on these reconstructed
product-law inputs.
% lean: CausalSmith.Stat.LdpOptvalueUniformFrontier.symmetricLaw_reconstructed_decisionLaw

\item \emph{Averaging through the sequential recursion.}
For \(0\le\ell\le n\), define the prefix message space
\[
\mathcal Z_{\le\ell}:=\prod_{i=1}^{\ell}\mathcal Z_i,
\qquad
\mathcal Z_{\le0}:=\{\langle\rangle\}.
\]
For a fixed observed-input vector
\(\mathbf o\in([d]\times\{0,1\}^2)^n\), paired-input vector
\(\mathbf v\in([d]\times\{-1,1\})^n\), and seed \(r\in\mathcal R\),
let \(\Pi_\ell^Q(\cdot\mid\mathbf o,r)\) and
\(\Pi_\ell^{\mathsf K}(\cdot\mid\mathbf v,r)\) be the prefix probability
laws defined by
\[
\Pi_0^Q=\Pi_0^{\mathsf K}=\delta_{\langle\rangle}
\]
and, for \(0\le\ell<n\),
\[
\begin{aligned}
\Pi_{\ell+1}^Q(dh,dz\mid\mathbf o,r)
&=\Pi_\ell^Q(dh\mid\mathbf o,r)\,
 Q_{\ell+1}(dz\mid o_{\ell+1},h,r),\\
\Pi_{\ell+1}^{\mathsf K}(dh,dz\mid\mathbf v,r)
&=\Pi_\ell^{\mathsf K}(dh\mid\mathbf v,r)\,
 \mathsf K_{\ell+1}(dz\mid v_{\ell+1},h,r),
\end{aligned}
\]
where \(h\in\mathcal Z_{\le\ell}\), \(z\in\mathcal Z_{\ell+1}\), and
\((h,z)\) denotes the appended history.

We prove by induction on \(\ell\) that, for every measurable
\(\varphi:\mathcal Z_{\le\ell}\to[0,\infty]\),
\[
\begin{aligned}
&2^{-n}\sum_{\mathbf a\in\{0,1\}^n}
 \int\varphi(h)\,
 \Pi_\ell^Q(dh\mid
 \operatorname{rec}^{(n)}(\mathbf v,\mathbf a),r)\\
&\qquad=\int\varphi(h)\,
 \Pi_\ell^{\mathsf K}(dh\mid\mathbf v,r).
\end{aligned}
\]
For \(\ell=0\), both prefix laws are concentrated on the empty history,
and \(2^{-n}\sum_{\mathbf a\in\{0,1\}^n}1=1\).

For the induction step, take measurable
\(\varphi:\mathcal Z_{\le\ell+1}\to[0,\infty]\) and define
\[
\varphi_{\mathrm{av},\ell}(h)
:=\int\varphi(h,z)\,
 \mathsf K_{\ell+1}(dz\mid v_{\ell+1},h,r),
\qquad h\in\mathcal Z_{\le\ell}.
\]
This is a measurable nonnegative function. For \(i\in[n]\) and
\(a\in\{0,1\}\), define the updated bit vector by
\[
(\mathbf a^{\,i\leftarrow a})_{i'}
:=
\begin{cases}
a,&i'=i,\\
a_{i'},&i'\ne i,
\end{cases}
\qquad i'\in[n].
\]
The uniform distribution on bit vectors is invariant under flipping
coordinate \(i\). Pairing each vector with its flipped vector shows that
averaging any nonnegative function over all bit vectors equals
averaging its two versions with coordinate \(i\) set to zero and one.

Apply this identity at \(i=\ell+1\). Updating that coordinate leaves the
first \(\ell\) inputs unchanged, hence leaves the preceding prefix law
unchanged. The recursion and the definition of \(\mathsf K_{\ell+1}\)
therefore give, for every \(\mathbf a\),
\[
\begin{aligned}
&\frac12\sum_{a\in\{0,1\}}
 \int\varphi(h,z)\,
 \Pi_{\ell+1}^Q
 \bigl(dh,dz\mid
 \operatorname{rec}^{(n)}
   (\mathbf v,\mathbf a^{\,\ell+1\leftarrow a}),r\bigr)\\
&\qquad=
 \int\left[
 \frac12\sum_{a\in\{0,1\}}
 \int\varphi(h,z)\,
 Q_{\ell+1}(dz\mid
 \operatorname{rec}(j_{\ell+1},s_{\ell+1},a),h,r)
 \right]\\
&\hspace{6em}\cdot
 \Pi_\ell^Q(dh\mid
 \operatorname{rec}^{(n)}(\mathbf v,\mathbf a),r)\\
&\qquad=
 \int\varphi_{\mathrm{av},\ell}(h)\,
 \Pi_\ell^Q(dh\mid
 \operatorname{rec}^{(n)}(\mathbf v,\mathbf a),r).
\end{aligned}
\]
Averaging this equality over \(\mathbf a\), applying the induction
hypothesis to \(\varphi_{\mathrm{av},\ell}\), and using the paired
recursion proves
\[
\begin{aligned}
&2^{-n}\sum_{\mathbf a}
 \int\varphi\,
 \Pi_{\ell+1}^Q
 (\,\cdot\mid\operatorname{rec}^{(n)}(\mathbf v,\mathbf a),r)\\
&\qquad=
 \int\varphi_{\mathrm{av},\ell}(h)\,
 \Pi_\ell^{\mathsf K}(dh\mid\mathbf v,r)
 =\int\varphi\,\Pi_{\ell+1}^{\mathsf K}
 (\,\cdot\mid\mathbf v,r).
\end{aligned}
\]
At \(\ell=n\), indicator functions of measurable transcript events give
the equality of probability laws
\[
2^{-n}\sum_{\mathbf a\in\{0,1\}^n}
\Pi_n^Q
(\,\cdot\mid\operatorname{rec}^{(n)}(\mathbf v,\mathbf a),r)
=\Pi_n^{\mathsf K}(\,\cdot\mid\mathbf v,r).
\]
Integrating against \(P_\theta^{\otimes n}\), \(\lambda\), and the uniform
analyst randomizer, and using the identity expressing
\(\mathcal L_{G_\theta}(\mathbf O,R,U)\) as the reconstruction
pushforward of \(\Lambda_\theta\), proves
\[
\begin{aligned}
&\mathcal L_{G_\theta,Q}(\mathbf Z,R,U)\\
&\qquad=
P_\theta^{\otimes n}\otimes\lambda
 \otimes\operatorname{Uniform}[0,1]
\ \text{followed by }\ 
\mathbf Z\sim\Pi_n^{\mathsf K}(\cdot\mid\mathbf v,r).
\end{aligned}
\]
Equivalently, for every measurable
\(\mathcal E\subseteq\mathcal Z_{\le n}\times\mathcal R\times[0,1]\),
the common probability is
\[
\int\!\int\!\int\!\int
\mathbf1_{\mathcal E}(\mathbf z,r,u)\,
\Pi_n^{\mathsf K}(d\mathbf z\mid\mathbf v,r)\,
\operatorname{Uniform}[0,1](du)\,
\lambda(dr)\,P_\theta^{\otimes n}(d\mathbf v).
\]
Projecting onto \((R,\mathbf Z)\) gives the first asserted joint
seed-transcript equality. The induction begins at zero and therefore
also covers \(n=0\).
% lean: CausalSmith.Stat.LdpOptvalueUniformFrontier.averagedProtocol_reconstructed_decisionLaw

\item \emph{Class preservation by signed preprocessing.}
For the converse direction, fix a paired-alphabet protocol \(\mathsf K\).
Define the deterministic paired-record map by
\[
\operatorname{pair}(j,a,y)
:=\bigl(j,(2a-1)(2y-1)\bigr),
\qquad
(j,a,y)\in[d]\times\{0,1\}^2,
\]
and define \(Q\), with the same seed law and message spaces, by
\[
Q_i(\cdot\mid(j,a,y),h,r)
:=\mathsf K_i(\cdot\mid\operatorname{pair}(j,a,y),h,r).
\]
Composition with this measurable finite-alphabet map preserves the
probability-kernel property.

If \(\mathsf K\) satisfies the paired privacy inequality, then for every
two observed records \(o,o'\), every history and seed, and every
measurable message event \(E\),
\[
\begin{aligned}
Q_i(E\mid o,h,r)
&=\mathsf K_i(E\mid\operatorname{pair}(o),h,r)\\
&\le e^\varepsilon
 \mathsf K_i(E\mid\operatorname{pair}(o'),h,r)
=e^\varepsilon Q_i(E\mid o',h,r).
\end{aligned}
\]
Thus \(Q\) satisfies \cref{obj:ass:full-record-privacy} and belongs to
the sequential class whenever \(\mathsf K\) does. If \(\mathsf K\) is
independent of history, then
\[
Q_i(\cdot\mid o,h,r)
=\mathsf K_i(\cdot\mid\operatorname{pair}(o),h,r)
=\mathsf K_i(\cdot\mid\operatorname{pair}(o),h',r)
=Q_i(\cdot\mid o,h',r).
\]
This supplies \cref{obj:ass:noninteraction} and hence the noninteractive
membership in \cref{obj:def:noninteractive-class}.
% lean: CausalSmith.Stat.LdpOptvalueUniformFrontier.signed_protocol_classes

\item \emph{The converse transcript identity.}
For a deterministic observed vector \(\mathbf o\), define its
coordinatewise paired image by
\[
\operatorname{pair}^{(n)}(\mathbf o)
:=\bigl(\operatorname{pair}(o_i)\bigr)_{i=1}^n.
\]
The previously established identities
\(\tau(G_\theta)=\theta\) and
\(\mathcal L_{G_\theta}(X,S)=P_\theta\) imply, by coordinatewise
pushforward of the independent observed sample,
\[
\begin{aligned}
&\bigl((G_\theta)_O^{\otimes n}
 \otimes\lambda\otimes\operatorname{Uniform}[0,1]\bigr)
 \circ
 \bigl[(\mathbf o,r,u)
 \mapsto(\operatorname{pair}^{(n)}(\mathbf o),r,u)\bigr]^{-1}\\
&\qquad=
P_\theta^{\otimes n}\otimes\lambda
 \otimes\operatorname{Uniform}[0,1].
\end{aligned}
\]
Here \cref{obj:ass:iid-people,obj:ass:independent-randomness} identify the
measure on the left before pushforward with the sampling scheme's
joint input law.

For each fixed \(\mathbf o,r\), induction through the prefix recursion
gives
\[
\Pi_\ell^Q(\cdot\mid\mathbf o,r)
=\Pi_\ell^{\mathsf K}
(\cdot\mid\operatorname{pair}^{(n)}(\mathbf o),r),
\qquad 0\le\ell\le n.
\]
The zero-length laws are the same point mass. At a successor step,
the induction hypothesis identifies the preceding prefix laws, and
the defining equality
\[
Q_{\ell+1}(dz\mid o_{\ell+1},h,r)
=\mathsf K_{\ell+1}
(dz\mid\operatorname{pair}(o_{\ell+1}),h,r)
\]
identifies their extension kernels. Thus their successor prefix laws
are equal as well. In particular, their full fixed-input transcript
laws agree at \(\ell=n\).

Integrating this fixed-input equality and using the displayed input
pushforward proves that the joint law of \((\mathbf Z,R,U)\) under
\(G_\theta,Q\) equals the law obtained from independent \(P_\theta\)
inputs, the seed law \(\lambda\), and \(\mathsf K\). Projecting onto
\((R,\mathbf Z)\) gives the asserted converse equality. The same
zero-length base case covers \(n=0\).
% lean: CausalSmith.Stat.LdpOptvalueUniformFrontier.pulledProtocol_decisionLaw
\end{enumerate}
\end{proof}

\begin{proof}[Proof of \cref{obj:thm:uniform-private-value-frontiers}]
We first establish uniform finite-resource bounds, then prove the numerical resource conclusions and transfer them to the minimax quantities. Throughout the finite-resource argument, fix an allowed tuple and a sampling scheme satisfying \cref{obj:ass:iid-people,obj:ass:independent-randomness}, and write
\[
t=n\varepsilon^2,\qquad L=\log(ed),\qquad
z_*=\frac{d^2L}{t},\qquad w=\log(e+z_*).
\]
Thus \(t>0\), and
\[
1\le L=1+\log d\le d.
\]

\begin{enumerate}
\item
Invoke the published Bernstein limit
\[
2k_{\mathrm{half}} e_{2k_{\mathrm{half}}}\longrightarrow\beta_*,
\qquad k_{\mathrm{half}}\in\mathbb N_{>0},
\qquad 0.278<\beta_*<0.286,
\]
as recorded in \citep[Section 3.1, p.~8]{CaiLow2011Nonsmooth}. Here
\[
e_K=\inf_{\deg p\le K}\sup_{|x|\le1}\bigl|p(x)-|x|\bigr|.
\]
Choose an integer \(N_{\mathrm{tail}}\) such that
\[
k_{\mathrm{half}}\ge N_{\mathrm{tail}}
\quad\Longrightarrow\quad
(2k_{\mathrm{half}})e_{2k_{\mathrm{half}}}\ge\frac1{10},
\]
and define
\[
N_{\mathrm B}=\max\{N_{\mathrm{tail}},1\},
\qquad
b_*=\min\left\{\frac1{10},e_{2N_{\mathrm B}}\right\}.
\]
The tail inequality gives \(e_{2N_{\mathrm B}}>0\), hence \(b_*>0\).
The approximation errors decrease with degree, because the admissible polynomial spaces increase. Therefore, for every positive even integer \(K\),
\[
e_K\ge\frac{b_*}{K}.
\]
Indeed, write \(K=2k_{\mathrm{half}}\). If \(k_{\mathrm{half}}\ge N_{\mathrm B}\), then \(Ke_K\ge1/10\ge b_*\). If \(k_{\mathrm{half}}<N_{\mathrm B}\), then
\[
Ke_K\ge e_K\ge e_{2N_{\mathrm B}}\ge b_*.
\]

For such a \(K\) and an amplitude \(0<\alpha\le1/2\), take the measures supplied by \cref{obj:lem:cai-low-moment-duality} and define
\[
\nu_\ell=(x\mapsto\alpha x)_{\#}\xi_\ell,
\qquad \ell\in\{0,1\}.
\]
These are probability measures supported on \([-\alpha,\alpha]\). Scaling their moments gives
\[
\int u^v\,\nu_0(du)=\int u^v\,\nu_1(du)
\quad(0\le v\le K),
\qquad
\int |u|\,\nu_1(du)-\int |u|\,\nu_0(du)
=2\alpha e_K.
\]
Define the coordinate priors, induced causal priors, welfare centers, and separation by
\[
\Pi_\ell=\nu_\ell^{\otimes d},\qquad
\Lambda_\ell=(\theta\mapsto G_\theta)_{\#}\Pi_\ell,
\qquad
v_\ell=\frac12+\frac12\int|u|\,\nu_\ell(du),
\qquad
\Delta=v_1-v_0=\alpha e_K.
\]
By \cref{obj:prop:signed-causal-bridge}, the priors \(\Lambda_\ell\) are supported on \(\mathcal V_d\), and
\[
V(G_\theta)=\frac12+\frac1{2d}\sum_{j=1}^d|\theta_j|.
\]
Independence of the prior coordinates consequently gives
\[
\mathbb E_{\Lambda_\ell}V(P)=v_\ell,
\qquad
\operatorname{Var}_{\Lambda_\ell}(V(P))
=\frac1{4d}\operatorname{Var}_{\nu_\ell}(|u|)
\le\frac{\alpha^2}{4d}.
\]
Here the last inequality follows from \(|u|\le\alpha\) almost surely. Chebyshev's inequality and the approximation lower bound yield
\[
\Delta\ge\frac{b_*\alpha}{K}>0,
\qquad
\Lambda_\ell\!\left\{|V(P)-v_\ell|>\frac{\Delta}{8}\right\}
\le\frac{16\alpha^2}{d\Delta^2}
\le\frac{16K^2}{b_*^2d}.
\]
% lean: CausalSmith.Stat.LdpOptvalueUniformFrontier.causal_dual_product_priors_target_concentration_of_gate

\item
Use the converse tuning in \cref{obj:def:frontier-handle}:
\[
(\alpha,k)=
\begin{cases}
\left(\dfrac{\sqrt{z_*}}{100},\,2\lceil16L\rceil+1\right),
&z_*\le1,\\[3pt]
\left(\dfrac1{100},\,2\lceil16L/w\rceil+1\right),
&z_*>1,
\end{cases}
\qquad K=k-1.
\]
Both branches have \(0<\alpha\le1/2\) and a positive even \(K\).
The ceiling inequalities give
\[
K\le
\begin{cases}
34L,&z_*\le1,\\
34\max\{1,L/w\},&z_*>1.
\end{cases}
\]
Since \(w\ge1\) and \(L\ge1\), both bounds imply \(K\le34L\).

We also have the calibrated separation inequality
\[
\frac{\alpha}{K}\ge\frac1{3400}\sqrt{\rho(t,d)}.
\]
For \(z_*\le1\), the dense formula gives
\[
\sqrt{\rho(t,d)}=\frac{\sqrt{z_*}}{L},
\qquad
\frac{\alpha}{K}
\ge\frac{\sqrt{z_*}}{3400L}.
\]
For \(z_*>1\), put
\[
u_{\mathrm{rate}}=\min\{1,w/L\}=\sqrt{\rho(t,d)}.
\]
Then \(u_{\mathrm{rate}}\le1\) and
\(u_{\mathrm{rate}}L/w\le1\). Hence
\[
u_{\mathrm{rate}}K
\le u_{\mathrm{rate}}(32L/w+2)\le34,
\qquad
\frac{\alpha}{K}\ge\frac{u_{\mathrm{rate}}}{3400}.
\]

Because
\[
\frac{L^2}{d}\longrightarrow0\qquad(d\longrightarrow\infty),
\]
choose a fixed integer \(D_{\mathrm{con}}\ge2\) such that
\[
d>D_{\mathrm{con}}
\quad\Longrightarrow\quad
\frac{16(34L)^2}{b_*^2d}\le\frac1{100}.
\]
For these dimensions, the priors constructed in the preceding step satisfy
\[
\Delta\ge\frac{b_*}{3400}\sqrt{\rho(t,d)},
\qquad
\Lambda_\ell\!\left\{|V(P)-v_\ell|>\frac{\Delta}{8}\right\}
\le\frac1{100}
\quad(\ell=0,1).
\]
The cutoff depends on the universal approximation constant and is uniform over the allowed resource tuples.
% lean: CausalSmith.Stat.LdpOptvalueUniformFrontier.causal_concentrated_resource_priors_of_gate

\item
Fix any \(Q\in\mathfrak Q_{\mathrm{SI}}(n,d,\varepsilon)\). The measurable-density requirement in \cref{obj:lem:adaptive-moment-contraction} is supplied by \cref{obj:lem:measurable-kernel-radon-nikodym}. Define
\[
\beta=\frac{e^\varepsilon-1}{2d},
\qquad
\gamma_{\mathrm{con}}
=\frac{\alpha\sqrt{et}}{d\sqrt{k}},
\qquad
w_{\mathrm{con}}=
\begin{cases}
1,&z_*\le1,\\
w,&z_*>1.
\end{cases}
\]
The matched moments through \(K=k-1\) give, when \(k\le n\),
\[
\operatorname{TV}(\mathsf M_0^Q,\mathsf M_1^Q)
\le d\alpha^k\sqrt{\binom nk}\,\beta^k
\le d\gamma_{\mathrm{con}}^k.
\]
To justify the second inequality, the budget range gives
\[
e^\varepsilon-1\le2\varepsilon,
\qquad \beta\le\frac{\varepsilon}{d}.
\]
Moreover, the exponential series implies \(k^k/k!\le e^k\), so
\[
\binom nk\le\frac{n^k}{k!}\le\left(\frac{en}{k}\right)^k,
\qquad
\sqrt{\binom nk}\le\left(\sqrt{\frac{en}{k}}\right)^k.
\]
For \(k>n\), \cref{obj:lem:adaptive-moment-contraction} instead gives
\(\mathsf M_0^Q=\mathsf M_1^Q\).

It remains to bound the geometric expression uniformly. In both branches,
\[
kw_{\mathrm{con}}\ge32L,
\qquad
\gamma_{\mathrm{con}}^2=\frac{\alpha^2eL}{kz_*}.
\]
For \(z_*\le1\), the amplitude formula and \(e<3\) give
\[
\alpha^2e\,w_{\mathrm{con}}e^{w_{\mathrm{con}}/2}
=\frac{z_*}{10000}e\,e^{1/2}
\le32z_*.
\]
For \(z_*>1\), use
\[
e^w=e+z_*,
\qquad
w\le2e^{w/2},
\]
the latter following from \(1+w/2\le e^{w/2}\). Thus
\[
we^{w/2}\le2(e+z_*)\le8z_*,
\qquad
\alpha^2e\,we^{w/2}\le32z_*.
\]
Combining either budget inequality with \(kw_{\mathrm{con}}\ge32L\) yields
\[
\gamma_{\mathrm{con}}^2
\le e^{-w_{\mathrm{con}}/2},
\qquad
\gamma_{\mathrm{con}}\le e^{-w_{\mathrm{con}}/4}.
\]
Since \(d\le e^L\) and \(L\ge1\),
\[
d\gamma_{\mathrm{con}}^k
\le de^{-kw_{\mathrm{con}}/4}
\le de^{-8L}
\le e^{-7L}
\le e^{-7}
\le\frac1{100}.
\]
For the last inequality, \(e\ge2\) gives \(e^7\ge128\).
Consequently, for every sequential protocol,
\[
\operatorname{TV}(\mathsf M_0^Q,\mathsf M_1^Q)\le\frac1{100}.
\]
The sampling and randomness assumptions identify these mixtures with the joint seed-transcript mixtures under the fixed sampling scheme; the independent analyst randomizer supplies the decision mixtures required by \cref{obj:lem:separated-value-mixtures}.
% lean: CausalSmith.Stat.LdpOptvalueUniformFrontier.causal_lawspace_resource_contraction

\item
Suppose first that \(d>D_{\mathrm{con}}\). Apply \cref{obj:lem:separated-value-mixtures} with
\[
\eta_0=\omega=\frac1{100}.
\]
The support, concentration, and transcript-proximity hypotheses were established above. For every estimator and every globally honest connected interval,
\[
\sup_{P\in\mathcal V_d}
\mathbb E_{P,Q}(T-V(P))^2
\ge\frac{9\Delta^2}{128}\frac{97}{100},
\]
\[
\sup_{P\in\mathcal V_d}
\mathbb E_{P,Q}\operatorname{length}I
\ge\frac{3\Delta}{4}\frac{77}{100}.
\]
Define
\[
s_{\mathrm{con}}=\frac{b_*}{3400},
\qquad
c_{\mathrm{large}}
=\min\left\{\frac{s_{\mathrm{con}}^2}{100},
                 \frac{s_{\mathrm{con}}}{100}\right\}>0.
\]
The separation bound gives
\[
\Delta^2\ge s_{\mathrm{con}}^2\rho(t,d),
\qquad
\Delta\ge s_{\mathrm{con}}\sqrt{\rho(t,d)}.
\]
Since \(9\cdot97/12800\ge1/100\) and
\(3\cdot77/400\ge1/100\), the preceding decision bounds imply
\[
\sup_{P\in\mathcal V_d}\mathbb E_{P,Q}(T-V(P))^2
\ge c_{\mathrm{large}}\rho(t,d),
\qquad
\sup_{P\in\mathcal V_d}\mathbb E_{P,Q}\operatorname{length}I
\ge c_{\mathrm{large}}\sqrt{\rho(t,d)}.
\]
% lean: CausalSmith.Stat.LdpOptvalueUniformFrontier.frontier_lower_of_gate

\item
For the bounded dimensions and the later resource analysis, we need
\[
\min\{1,t^{-1}\}\le\rho(t,d)
\le4\min\{1,d^2/t\},
\qquad
\frac14\min\{t,d\}\le t\rho(t,d).
\]
Here is their branchwise derivation. In the dense branch,
\[
\rho(t,d)=\frac{d^2}{tL},
\qquad
\frac1t\le\rho(t,d)\le\min\{1,d^2/t\},
\qquad
t\rho(t,d)=\frac{d^2}{L}\ge d,
\]
using \(1\le L\le d\).

For the nondense branch, \(z_*>1\) and
\[
\rho(t,d)=\min\{1,(w/L)^2\}.
\]
The inequalities \(\log d\ge1-d^{-1}\) and \(e<3\) give
\[
dL=d(1+\log d)\ge2d-1\ge3>e.
\]
Consequently,
\[
\log(e+d^2L)\ge L,
\qquad
\log(e+dL)\ge\frac L2.
\]
For the second inequality, the logarithm on the left is at least both
\(1\) and \(\log d\).

The function
\[
\Phi_{\mathrm{res}}(t')
=t'\left[\log\left(e+\frac{d^2L}{t'}\right)\right]^2,
\qquad t'\in(0,\infty),
\]
is increasing. Indeed, for every \(v_{\mathrm{shift}}\ge0\),
\[
\log(e+v_{\mathrm{shift}})
\ge2-\frac e{e+v_{\mathrm{shift}}},
\]
by \(\log x\ge1-x^{-1}\) applied to
\(x=(e+v_{\mathrm{shift}})/e\). Therefore
\[
\log(e+v_{\mathrm{shift}})
-\frac{2v_{\mathrm{shift}}}{e+v_{\mathrm{shift}}}
\ge\frac e{e+v_{\mathrm{shift}}}\ge0,
\]
and differentiation gives
\[
\Phi_{\mathrm{res}}'(t')
=
\log\left(e+\frac{d^2L}{t'}\right)
\left[
\log\left(e+\frac{d^2L}{t'}\right)
-\frac{2d^2L/t'}{e+d^2L/t'}
\right]\ge0.
\]

If \(t\le1\), then \(w\ge\log(e+d^2L)\ge L\), so \(\rho(t,d)=1\).
If \(t\ge1\), monotonicity gives
\[
tw^2\ge[\log(e+d^2L)]^2\ge L^2,
\]
which proves \(\rho(t,d)\ge\min\{1,t^{-1}\}\).
If \(t\le d\), then \(w\ge\log(e+dL)\ge L/2\), so
\[
t(w/L)^2\ge t/4.
\]
If \(t\ge d\), monotonicity instead gives
\[
tw^2\ge d[\log(e+dL)]^2\ge dL^2/4.
\]
Since \(t\ge\min\{t,d\}/4\), taking the minimum with \(t\) proves the asserted lower bound for \(t\rho\).

Finally, for \(z_*\ge1\),
\[
\log(e+z_*)
\le\log(e+1)+\log z_*
\le2+2(\sqrt{z_*}-1)=2\sqrt{z_*}.
\]
Here \(\log(e+1)\le2\), and
\(\log\sqrt{z_*}\le\sqrt{z_*}-1\).
Hence
\[
(w/L)^2\le\frac{4z_*}{L^2}
=\frac{4d^2}{tL}\le\frac{4d^2}{t}.
\]
Together with \(\rho\le1\), this proves the upper comparison.
% lean: CausalSmith.Stat.LdpOptvalueUniformFrontier.rho_elementary_comparisons

\item
We next obtain a per-decision two-point bound for every dimension. Define
\[
\alpha_{\mathrm{pt}}
=\frac18\min\{1,t^{-1/2}\},
\qquad
\theta^{(0)}_j=0,\qquad
\theta^{(1)}_j=\alpha_{\mathrm{pt}}
\quad(j\in[d]).
\]
These vectors belong to \([-1/2,1/2]^d\), and
\[
0<\alpha_{\mathrm{pt}}\le\frac18,\qquad
\alpha_{\mathrm{pt}}\sqrt t\le\frac18.
\]
By \cref{obj:prop:signed-causal-bridge},
\[
G_{\theta^{(0)}},G_{\theta^{(1)}}\in\mathcal V_d,
\qquad
V(G_{\theta^{(0)}})=\frac12,\qquad
V(G_{\theta^{(1)}})=\frac12+\frac{\alpha_{\mathrm{pt}}}{2}.
\]
Apply \cref{obj:lem:adaptive-moment-contraction} with matching order \(1\) and coordinate measures \(\delta_0,\delta_{\alpha_{\mathrm{pt}}}\); their zeroth moments agree. For their seed-transcript laws,
\[
\operatorname{TV}
\le d\alpha_{\mathrm{pt}}\sqrt n\,
       \frac{e^\varepsilon-1}{2d}
\le\alpha_{\mathrm{pt}}\sqrt t
\le\frac18.
\]
The point priors have zero target-concentration error. Thus \cref{obj:lem:separated-value-mixtures}, with centers separated by
\(\alpha_{\mathrm{pt}}/2\), gives
\[
\sup_{P\in\mathcal V_d}\mathbb E_{P,Q}(T-V(P))^2
\ge\frac{63\alpha_{\mathrm{pt}}^2}{4096},
\qquad
\sup_{P\in\mathcal V_d}\mathbb E_{P,Q}\operatorname{length}I
\ge\frac{81\alpha_{\mathrm{pt}}}{320}.
\]
Since
\[
\alpha_{\mathrm{pt}}^2=\frac1{64}\min\{1,t^{-1}\},
\]
these imply
\[
\sup_{P\in\mathcal V_d}\mathbb E_{P,Q}(T-V(P))^2
\ge\frac1{8192}\min\{1,t^{-1}\},
\qquad
\sup_{P\in\mathcal V_d}\mathbb E_{P,Q}\operatorname{length}I
\ge\frac1{8192}\min\{1,t^{-1/2}\}.
\]

For \(d\le D_{\mathrm{con}}\), the elementary upper comparison gives
\[
\rho(t,d)\le4D_{\mathrm{con}}^2\min\{1,t^{-1}\}.
\]
To see the comparison inside the minimum, use
\(\min\{1,d^2/t\}\le1\) when \(t\le1\), and
\(\min\{1,d^2/t\}\le d^2/t\) when \(t>1\).
Taking square roots and enlarging the factor gives
\[
\sqrt{\rho(t,d)}
\le2D_{\mathrm{con}}\min\{1,t^{-1/2}\}
\le4D_{\mathrm{con}}^2\min\{1,t^{-1/2}\}.
\]
Define
\[
c_{\mathrm{small}}=\frac1{32768D_{\mathrm{con}}^2},
\qquad
c_{\mathrm{lo}}=\min\{c_{\mathrm{large}},c_{\mathrm{small}}\}>0.
\]
For \(d\le D_{\mathrm{con}}\), the two-point bounds give the frontier lower bounds with \(c_{\mathrm{small}}\); for \(d>D_{\mathrm{con}}\), the product-prior bounds give them with \(c_{\mathrm{large}}\). Thus, for every sequential protocol and every admissible decision,
\[
\sup_{P\in\mathcal V_d}\mathbb E_{P,Q}(T-V(P))^2
\ge c_{\mathrm{lo}}\rho(t,d),
\qquad
\sup_{P\in\mathcal V_d}\mathbb E_{P,Q}\operatorname{length}I
\ge c_{\mathrm{lo}}\sqrt{\rho(t,d)}.
\]
Every noninteractive protocol is sequential by
\cref{obj:def:noninteractive-class,obj:def:sequential-class}.
Taking the decision and protocol infima therefore gives, for either class,
\[
\mathfrak R_{\mathcal C}\ge c_{\mathrm{lo}}\rho(t,d),
\qquad
\mathfrak H_{\mathcal C}\ge c_{\mathrm{lo}}\sqrt{\rho(t,d)}.
\]
The risk lower bounds remain valid for extended nonnegative expectations.
% lean: CausalSmith.Stat.LdpOptvalueUniformFrontier.frontier_minimax_lower_of_gate

\item
For the upper bounds, use the construction \(\mathscr H\) of
\cref{obj:def:frontier-handle}, with its singleton public seed and transcript-based output \(\widehat V\). Its channels are fixed by the public resources. For \(n=2\) the messages are constant. For \(n>2\), use the baseline, pilot, and evaluation blocks specified there, with
\[
m=\lfloor n/3\rfloor,\qquad m_0=n-2m,\qquad
\delta=\tanh(\varepsilon/2)=\frac{e^\varepsilon-1}{e^\varepsilon+1},
\qquad b=\frac d\delta,\qquad \sigma^2=\frac{b^2}{m}.
\]
The block sizes satisfy
\[
m\ge n/5>0,\qquad m_0\ge n/3>0,\qquad 2m\le n.
\]
For the first inequality, write \(n=3m+r\), \(r\in\{0,1,2\}\), and use \(m\ge1\), giving \(n\le5m\). The second follows from \(m\le n/3\).

The binary channel probabilities are between \((1-\delta)/2\) and
\((1+\delta)/2\), and
\[
\frac{1+\delta}{1-\delta}=e^\varepsilon.
\]
Thus their ratio across any two original records is at most
\(e^\varepsilon\). The vector channel satisfies
\[
2^{-d}(1-\delta)
\le Q^{\mathrm{vec}}(z\mid O=(j,A,Y))
\le 2^{-d}(1+\delta)
\qquad\bigl(z\in\{-1,1\}^d\bigr).
\]
Its probability normalization and noninteractive privacy are also supplied by
\cref{obj:lem:private-moment-and-dependence}. Summing the pointwise privacy inequalities over message sets proves full-record privacy. The kernels are independent of message history, so the entire construction belongs to the noninteractive class.

The sampling and randomness assumptions imply the product transcript law
\[
\mathcal L_P(\mathbf Z)
=\bigotimes_{i=1}^n
\left[\int Q_i(\,\cdot\mid o)\,P_O(do)\right].
\]
In particular, the disjoint blocks are independent and each vector block has the law used in \cref{obj:lem:finite-private-polynomial-calibration}.

The privacy attenuation obeys
\[
\delta\ge\varepsilon/4>0.
\]
Indeed, \(e^\varepsilon\ge1+\varepsilon\) and \(e^\varepsilon<3\) imply
\(4(e^\varepsilon-1)\ge\varepsilon(e^\varepsilon+1)\).
Consequently,
\[
\sigma^2=\frac{d^2}{m\delta^2}\le\frac{80d^2}{t}.
\]
For the baseline messages,
\[
\mathbb E_P\mathsf H_i=\delta(2B(P)-1),
\qquad
\operatorname{Var}_P(\mathsf H_i)\le1.
\]
Independence therefore gives
\[
\mathbb E_P\widehat B=B(P),
\qquad
\mathbb E_P(\widehat B-B(P))^2
\le\frac1{4m_0\delta^2}\le\frac{12}{t}.
\]

The interior arm means imply
\[
V(P)=\frac1d\sum_{j=1}^d
       \max\{\mu_{0j}(P),\mu_{1j}(P)\}\in[1/4,3/4].
\]
For \(n>2\), the decomposition in \cref{obj:prop:signed-causal-bridge} and clipping to this interval give
\[
(\widehat V-V(P))^2
\le2(\widehat B-B(P))^2
   +\frac12(\widehat F-F(\tau(P)))^2.
\]
This follows by first contracting the distance through clipping and then applying
\((x+y/2)^2\le2x^2+y^2/2\).
Hence
\[
\mathbb E_P(\widehat V-V(P))^2
\le\frac{24}{t}
+\frac12\mathbb E_P(\widehat F-F(\tau(P)))^2.
\]
All branches also satisfy the pointwise bound
\[
(\widehat V-V(P))^2\le\frac14.
\]
% lean: CausalSmith.Stat.LdpOptvalueUniformFrontier.frontier_product_risk_le_components

\item
We prove the risk bound with the fixed constant
\[
c_{\mathrm{up}}=10^{16}
\]
in the following order.

If \(n=2\), then \(\widehat V=1/2\) and \(t\le2\). Therefore
\[
(\widehat V-V(P))^2\le\frac1{16},
\qquad
\rho(t,d)\ge\min\{1,t^{-1}\}\ge\frac12,
\]
which proves the bound. For \(n>2\), if \(\rho(t,d)=1\), clipping gives
\[
\mathbb E_P(\widehat V-V(P))^2\le\frac14
\le c_{\mathrm{up}}\rho(t,d).
\]

Next suppose the selected branch is the fallback \(\widehat F=1/4\).
This branch has \(t<d^2L\) and, with
\[
H=\log(e+\sigma^2L),
\]
it satisfies \(L/(1024H)<4\). The noise bound gives
\[
H\le\log\!\bigl(80(e+z_*)\bigr)
=\log80+w\le6w,
\]
because \(\log80\le5\) and \(w\ge1\). Thus
\[
L<4096H\le24576w,
\qquad
\rho(t,d)=\bigl(\min\{1,w/L\}\bigr)^2
\ge\frac1{24576^2}.
\]
Clipping again yields
\[
\mathbb E_P(\widehat V-V(P))^2
\le\frac14
\le c_{\mathrm{up}}\rho(t,d).
\]

Next, when \(L<4096\) and \(t<d^2L\), use \(w\ge1\) to obtain
\[
\rho(t,d)\ge\frac1{4096^2}.
\]
The same clipping argument proves the bound.

If \(L<4096\) in the remaining case, then \(t\ge d^2L\), and the construction uses the absolute evaluation means. The private-message moments in
\cref{obj:lem:private-moment-and-dependence} imply
\[
\mathbb E_P\overline W_j=\tau_j(P),
\qquad
\mathbb E_P(\overline W_j-\tau_j(P))^2\le\sigma^2.
\]
The absolute-value inequality and the squared-average inequality give
\[
\left[
\frac1d\sum_{j=1}^d|\overline W_j|
-\frac1d\sum_{j=1}^d|\tau_j(P)|
\right]^2
\le\frac1d\sum_{j=1}^d
       (\overline W_j-\tau_j(P))^2.
\]
Hence
\[
\mathbb E_P(\widehat F-F(\tau(P)))^2\le\sigma^2
\]
and
\[
\mathbb E_P(\widehat V-V(P))^2
\le\frac{24}{t}+\frac{\sigma^2}{2}
\le\frac{64d^2}{t}
=64L\rho(t,d)
\le64\cdot4096\,\rho(t,d)
\le c_{\mathrm{up}}\rho(t,d).
\]
% lean: CausalSmith.Stat.LdpOptvalueUniformFrontier.frontier_upper_of_gate

\item
In the remaining dense case, \(L\ge4096\) and \(t\ge d^2L\).
The construction uses
\[
D=2\left\lfloor\frac{L}{2048}\right\rfloor,
\qquad T_*=4\sigma\sqrt L,\qquad
a_{\mathrm{rad}}=8\sigma\sqrt L.
\]
This is an even degree, and its moments fit the evaluation block:
\[
2D\le\frac{L}{512}
\le\frac{d^2L}{5}
\le\frac n5\le m.
\]
Here \(t\le n\) follows from \(\varepsilon\le1\).

The concentration and approximation hypotheses of
\cref{obj:lem:finite-private-polynomial-calibration} are supplied by
\cref{obj:lem:bounded-mean-concentration,obj:lem:cai-low-chebyshev-approximation}.
Its hybrid conclusion, applied to the independent pilot and evaluation blocks, gives
\[
\mathbb E_P(\widehat F-F(\tau(P)))^2
\le500000000\,\frac{\sigma^2}{L}
\le40000000000\,\rho(t,d).
\]
Since \(L\le d^2\), the dense formula also gives \(t^{-1}\le\rho(t,d)\).
Therefore
\[
\mathbb E_P(\widehat V-V(P))^2
\le\left(24+\frac{40000000000}{2}\right)\rho(t,d)
\le c_{\mathrm{up}}\rho(t,d).
\]
% lean: CausalSmith.Stat.LdpOptvalueUniformFrontier.frontierEstimator_hybrid_dense_rate

\item
The final case has \(L\ge4096\), \(t<d^2L\), \(\rho(t,d)\ne1\), and a branch distinct from the fallback. With \(H=\log(e+\sigma^2L)\), this means
\[
\frac{L}{1024H}\ge4,
\qquad
D=2\left\lfloor\frac{L}{2048H}\right\rfloor,
\qquad a_{\mathrm{rad}}=\frac12.
\]
Since \(H\ge1\), the floor inequalities give
\[
D\ge2,\qquad D\ \text{even},\qquad
\frac{L}{2048H}\le D\le\frac{L}{1024H}.
\]

We verify \(2D\le m\). Since \(0<\delta\le1\),
\[
\sigma^2\ge\frac{d^2}{m}.
\]
If \(m<2D\), the degree budget would imply
\[
m<2D\le\frac{L}{512H}\le\frac{L}{512},
\qquad
\sigma^2L>512d^2.
\]
Thus \(e+\sigma^2L\ge ed\) and \(H\ge L\). But then
\(D\le L/(1024H)\le1/1024\), contradicting \(D\ge2\).
This proves the moment-fit condition.

The global-polynomial conclusion of
\cref{obj:lem:finite-private-polynomial-calibration} now gives
\[
\mathbb E_P(\widehat F-F(\tau(P)))^2
\le\frac{e^{9DH}}{2d}+\frac1{(D+1)^2}.
\]
For completeness, the exponential dimension estimate used to calibrate this bound is
\[
L^2e^{L/100}\le4d.
\]
Indeed, \(L/4\le e^{L/4}\) gives \(L^2\le16e^{L/2}\), while
\(L\ge4096\) gives \(4\le e^{49L/100-1}\). Therefore
\[
L^2e^{L/100}
\le16e^{51L/100}
\le4e^{L-1}=4d.
\]
Since \(DH\le L/1024\), it follows that
\[
\frac{e^{9DH}}{2d}
\le\frac{e^{L/100}}{2d}\le\frac2{L^2}.
\]
As above, \(H\le6w\). The lower rounded-degree bound gives
\[
\frac1{D+1}\le\frac{2048H}{L}
\le12288\,\frac wL.
\]
The nonsaturation condition and the nondense formula imply
\[
\frac wL<1,\qquad \rho(t,d)=(w/L)^2.
\]
Using \(w\ge1\), we obtain
\[
\frac{e^{9DH}}{2d}\le2\rho(t,d),
\qquad
\frac1{(D+1)^2}\le150994944\,\rho(t,d).
\]
Also \(t^{-1}\le\rho(t,d)\): otherwise
\(\min\{1,t^{-1}\}=1\), and the elementary lower bound together with
\(\rho\le1\) would force saturation. Consequently,
\[
\mathbb E_P(\widehat V-V(P))^2
\le\left[24+\frac12(2+150994944)\right]\rho(t,d)
\le c_{\mathrm{up}}\rho(t,d).
\]
All cases have therefore established
\[
\sup_{P\in\mathcal V_d}
\mathbb E_P(\widehat V-V(P))^2
\le c_{\mathrm{up}}\rho(t,d).
\]
% lean: CausalSmith.Stat.LdpOptvalueUniformFrontier.frontier_global_resource_bound

\item
Use the interval in \cref{obj:def:frontier-handle}:
\[
q_*=\sqrt{10^{17}\rho(t,d)},
\qquad
I=
\left[
\max\{1/4,\widehat V-q_*\},
\min\{3/4,\widehat V+q_*\}
\right].
\]
The rate is positive by the elementary lower comparison, so \(q_*>0\).
Because \(\widehat V\in[1/4,3/4]\), these endpoints are ordered, and the interval is connected and closed. All decisions are Borel functions of the finite transcript.

For \(V(P)\in[1/4,3/4]\), failure of coverage implies
\(|\widehat V-V(P)|>q_*\). Markov's inequality therefore gives
\[
\Pr_P\{V(P)\notin I\}
\le\Pr_P\{(\widehat V-V(P))^2\ge q_*^2\}
\le\frac{10^{16}\rho(t,d)}{10^{17}\rho(t,d)}
=\frac1{10}.
\]
Clipping the endpoints gives
\[
\operatorname{length}I\le2q_*
=\sqrt{40c_{\mathrm{up}}}\,\sqrt{\rho(t,d)}.
\]
Thus the same noninteractive protocol, transcript-only estimator, and interval simultaneously satisfy
\[
\mathbb E_P(\widehat V-V(P))^2\le c_{\mathrm{up}}\rho(t,d),
\qquad
\Pr_P\{V(P)\in I\}\ge0.90,
\qquad
\mathbb E_P\operatorname{length}I
\le\sqrt{40c_{\mathrm{up}}}\,\sqrt{\rho(t,d)}.
\]
These hold for every \(P\in\mathcal V_d\), including the \(n=2\) branch.
% lean: CausalSmith.Stat.LdpOptvalueUniformFrontier.frontier_upper_of_gate

\item
Choose the common constants
\[
c=\min\{c_{\mathrm{lo}},1\},
\qquad
C_0=\max\{1,c_{\mathrm{up}},\sqrt{40c_{\mathrm{up}}}\}.
\]
Then
\[
0<c\le C_0<\infty,\qquad
c\le c_{\mathrm{lo}},\qquad
c_{\mathrm{up}}\le C_0,\qquad
\sqrt{40c_{\mathrm{up}}}\le C_0.
\]
The displayed attainment bounds consequently hold with \(C_0\).
The estimate is a Borel function of \(\mathbf Z\) alone, since the public seed is a singleton and the construction ignores analyst randomization.

The attaining protocol is admissible for both classes. Using it and its estimate as candidates in \cref{obj:def:risk}, and its globally honest interval as a candidate in \cref{obj:def:honest-length}, gives
\[
\mathfrak R_{\mathcal C}
\le\sup_{P\in\mathcal V_d}\mathbb E_P(\widehat V-V(P))^2
\le C_0\rho(t,d),
\]
\[
\mathfrak H_{\mathcal C}
\le\sup_{P\in\mathcal V_d}\mathbb E_P\operatorname{length}I
\le C_0\sqrt{\rho(t,d)}.
\]
Together with the lower bounds, this proves
\[
c\rho(t,d)\le\mathfrak R_{\mathcal C}\le C_0\rho(t,d),
\qquad
c\sqrt{\rho(t,d)}
\le\mathfrak H_{\mathcal C}\le C_0\sqrt{\rho(t,d)}.
\]
The attaining candidates are precisely the resource-selected construction, so its three guarantees hold with this same constant.
% lean: CausalSmith.Stat.LdpOptvalueUniformFrontier.causalAttainment_minimax_upper

\item
Protocol-class inclusion gives
\[
\mathfrak R_{\mathrm{SI}}\le\mathfrak R_{\mathrm{NI}},
\qquad
\mathfrak H_{\mathrm{SI}}\le\mathfrak H_{\mathrm{NI}}.
\]
Both denominators are positive by the lower bounds, and all four quantities are finite by the upper bounds. Hence
\[
1\le\frac{\mathfrak R_{\mathrm{NI}}}{\mathfrak R_{\mathrm{SI}}}
\le\frac{C_0\rho(t,d)}{c\rho(t,d)}
=\frac{C_0}{c},
\qquad
1\le\frac{\mathfrak H_{\mathrm{NI}}}{\mathfrak H_{\mathrm{SI}}}
\le\frac{C_0\sqrt{\rho(t,d)}}{c\sqrt{\rho(t,d)}}
=\frac{C_0}{c}.
\]
% lean: CausalSmith.Stat.LdpOptvalueUniformFrontier.frontier_ratio_comparison

\item
Squaring the nonnegative length bounds gives
\[
c^2\rho(t,d)\le\mathfrak H_{\mathcal C}^2
\le C_0^2\rho(t,d).
\]
The risk bounds then yield the two comparison chains
\[
\frac{c^2}{C_0}\mathfrak R_{\mathcal C}
\le\frac{c^2}{C_0}\,C_0\rho(t,d)
=c^2\rho(t,d)
\le\mathfrak H_{\mathcal C}^2,
\]
\[
\mathfrak H_{\mathcal C}^2
\le C_0^2\rho(t,d)
=\frac{C_0^2}{c}\,c\rho(t,d)
\le\frac{C_0^2}{c}\mathfrak R_{\mathcal C}.
\]
% lean: CausalSmith.Stat.LdpOptvalueUniformFrontier.frontier_square_comparison

\item
We now establish the resource conclusions. For any \(d\ge2\) and
\(0<t<d^2L\), the nondense expression has a positive logarithmic ratio, and therefore
\[
\begin{aligned}
\rho(t,d)=1
&\iff \log(e+d^2L/t)\ge L\\
&\iff e+d^2L/t\ge ed\\
&\iff t\le\frac{d^2L}{ed-e}.
\end{aligned}
\]
The final equivalence uses \(t>0\) and \(ed-e>0\).
% lean: CausalSmith.Stat.LdpOptvalueUniformFrontier.rho_saturation_iff

\item
Fix any sequence of allowed tuples and define
\[
t_k=n_k\varepsilon_k^2,\qquad
L_k=\log(ed_k),\qquad
\rho_k=\rho(t_k,d_k),\qquad
g_k=\frac{\log(1+d_k^2/t_k)}{L_k}.
\]
All \(t_k\) are positive, and \(\rho_k,g_k\ge0\).

Two numerical inequalities will give the consistency equivalence. For arbitrary \(d\ge2\), \(t>0\), define
\[
g(t,d)=\frac{\log(1+d^2/t)}{L},
\qquad
e_{\mathrm{corr}}(d)=\frac{\log(e+L)}{L}.
\]
For \(x_{\mathrm{res}}=d^2/t\ge0\), the inequalities
\[
1+x_{\mathrm{res}}\le e+Lx_{\mathrm{res}}
\le(1+x_{\mathrm{res}})(e+L)
\]
imply
\[
\log(1+x_{\mathrm{res}})
\le\log(e+Lx_{\mathrm{res}})
\le\log(1+x_{\mathrm{res}})+\log(e+L).
\]
It follows that
\[
\rho(t,d)<1\quad\Longrightarrow\quad
g(t,d)^2\le\rho(t,d).
\]
In the nondense branch, this follows by dividing the first logarithmic inequality by \(L\), because nonsaturation selects the untruncated ratio. In the dense branch,
\[
x_{\mathrm{res}}L\le1,\qquad
0\le g(t,d)\le\frac{x_{\mathrm{res}}}{L}\le1,
\]
using \(\log(1+x_{\mathrm{res}})\le x_{\mathrm{res}}\); hence
\[
g(t,d)^2\le\frac{x_{\mathrm{res}}}{L}=\rho(t,d).
\]
For both branches we also have
\[
\rho(t,d)
\le\bigl[g(t,d)+e_{\mathrm{corr}}(d)\bigr]^2+\frac1{L^2}.
\]
In the nondense branch use the upper logarithmic inequality and the nonnegative truncated ratio. In the dense branch use
\(\rho(t,d)\le L^{-2}\).
Finally,
\[
e_{\mathrm{corr}}(d)\longrightarrow0,
\qquad
L^{-1}\longrightarrow0
\qquad(d\longrightarrow\infty),
\]
by \(L\to\infty\) and \(\log(e+L)/L\to0\).

Suppose \(\rho_k\to0\). The elementary lower comparison forces \(t_k\to\infty\): for any real threshold \(M_{\mathrm{test}}\), if
\(t_k<M_{\mathrm{test}}\), then
\[
\rho_k\ge
\min\{1,t_k^{-1}\}
\ge\frac1{\max\{M_{\mathrm{test}},1\}}>0.
\]
Such indices are eventually excluded by convergence to zero.
Eventually \(\rho_k<1\), and the first numerical inequality gives
\[
0\le g_k\le\sqrt{\rho_k}\longrightarrow0.
\]

Conversely, suppose \(t_k\to\infty\) and \(g_k\to0\). Fix a tolerance
\(\zeta>0\), and define
\[
q_{\mathrm{cons}}=\frac{\sqrt{\zeta}}4.
\]
Choose an integer \(D_{\mathrm{cons}}\) such that, for every
\(d\ge D_{\mathrm{cons}}\),
\[
e_{\mathrm{corr}}(d)<q_{\mathrm{cons}},
\qquad
\frac1{\log(ed)}<q_{\mathrm{cons}}.
\]
Eventually,
\[
g_k<q_{\mathrm{cons}},
\qquad
t_k\ge\frac{4D_{\mathrm{cons}}^2}{\zeta}+1.
\]
If \(d_k\ge D_{\mathrm{cons}}\), the numerical upper bound gives
\[
\rho_k<(2q_{\mathrm{cons}})^2+q_{\mathrm{cons}}^2
=\frac{5\zeta}{16}<\zeta.
\]
If \(d_k<D_{\mathrm{cons}}\), the elementary comparison gives
\[
\rho_k\le\frac{4d_k^2}{t_k}
\le\frac{4D_{\mathrm{cons}}^2}{t_k}<\zeta.
\]
Thus \(\rho_k\to0\), establishing the criterion for arbitrary allowed sequences.
% lean: CausalSmith.Stat.LdpOptvalueUniformFrontier.rho_consistency_iff

\item
We next prove the private-parametric criterion. Write eventual comparability as the existence of constants
\[
0<c_{\mathrm{cmp}}\le C_{\mathrm{cmp}}<\infty,
\qquad
\frac{c_{\mathrm{cmp}}}{t_k}\le\rho_k
\le\frac{C_{\mathrm{cmp}}}{t_k}
\quad\text{eventually}.
\]
If these inequalities hold, then \(\rho_k\le1\) implies
\(t_k\ge c_{\mathrm{cmp}}\) eventually. The elementary scaled lower bound also gives
\[
\frac14\min\{t_k,d_k\}\le t_k\rho_k\le C_{\mathrm{cmp}},
\]
so \(\min\{t_k,d_k\}\le4C_{\mathrm{cmp}}\) eventually.

Conversely, suppose that for some \(a>0\) and \(M\in\mathbb R\),
\[
t_k\ge a,\qquad \min\{t_k,d_k\}\le M
\quad\text{eventually}.
\]
Define
\[
M_+=\max\{1,M\},\qquad
c_{\mathrm{cmp}}=\min\{1,a\},\qquad
C_{\mathrm{cmp}}=4M_+^2.
\]
The elementary lower comparison gives
\[
\frac{c_{\mathrm{cmp}}}{t_k}
\le\min\{1,t_k^{-1}\}\le\rho_k.
\]
For the upper bound, if \(t_k\le d_k\), then \(t_k\le M_+\), and
\[
\rho_k\le1\le\frac{4M_+^2}{t_k}.
\]
If \(d_k\le t_k\), then \(d_k\le M_+\), and
\[
\rho_k\le\frac{4d_k^2}{t_k}
\le\frac{4M_+^2}{t_k}.
\]
This proves the equivalence.

When \(t_k\to\infty\), an eventual positive lower bound on \(t_k\) is automatic. If \(\min\{t_k,d_k\}\le M\) eventually, then eventually \(t_k>M\), forcing \(d_k\le M\). Enlarging \(M\) to an integer yields the stated eventual dimension bound. Conversely, an eventual integer bound on \(d_k\) bounds the minimum.
% lean: CausalSmith.Stat.LdpOptvalueUniformFrontier.rho_parametric_comparable_iff

\item
Fix resource-ratio endpoints
\[
0<a\le b_{\mathrm{ratio}}<\infty.
\]
Choose \(D_0\) sufficiently large that, for every \(d\ge D_0\), with \(L=\log(ed)\),
\[
L\ge b_{\mathrm{ratio}}+1,\qquad
2\log b_{\mathrm{ratio}}\le\log L,\qquad
\log(e+1/a)\le\log L.
\]
These thresholds exist because \(L,\log L\to\infty\).
For \(d\ge\max\{D_0,2\}\) and \(a\le t/d^2\le b_{\mathrm{ratio}}\), define
\[
u_{\mathrm{elbow}}=\frac{t}{d^2}\in[a,b_{\mathrm{ratio}}].
\]
Then \(t<d^2L\), and
\[
\log(e+L/u_{\mathrm{elbow}})
\ge\log(L/b_{\mathrm{ratio}})
=\log L-\log b_{\mathrm{ratio}}\ge\frac12\log L.
\]
Also,
\[
e+\frac{L}{u_{\mathrm{elbow}}}
\le L(e+1/a),
\]
because \(L\ge1\). Hence
\[
\log(e+L/u_{\mathrm{elbow}})
\le\log L+\log(e+1/a)\le2\log L.
\]
Since \(0\le\log L/L\le1\), truncation at one preserves the bounds
\[
\frac12\frac{\log L}{L}
\le\min\left\{1,\frac{\log(e+L/u_{\mathrm{elbow}})}L\right\}
\le2\frac{\log L}{L}.
\]
Squaring gives the uniform comparison
\[
\frac14\left(\frac{\log L}{L}\right)^2
\le\rho(t,d)\le
4\left(\frac{\log L}{L}\right)^2.
\]
This proves the final numerical resource assertion, with constants \(1/4\) and \(4\).
% lean: CausalSmith.Stat.LdpOptvalueUniformFrontier.rate_resource_conclusions

\item
Finally, fix a family of sampling schemes satisfying the two stated assumptions at each \(n,d\), and fix a protocol class \(\mathcal C\).
Along an allowed sequence define
\[
\mathfrak R_k=\mathfrak R_{\mathcal C}(n_k,d_k,\varepsilon_k),
\qquad
\mathfrak H_k=\mathfrak H_{\mathcal C}(n_k,d_k,\varepsilon_k).
\]
The finite upper bounds make both quantities finite, and their real-valued bounds are
\[
c\rho_k\le\mathfrak R_k\le C_0\rho_k,
\qquad
c\sqrt{\rho_k}\le\mathfrak H_k\le C_0\sqrt{\rho_k}.
\]
The first sandwich gives
\[
\mathfrak R_k\to0\iff\rho_k\to0.
\]
The second, together with continuity of the square root and
\((\sqrt{\rho_k})^2=\rho_k\), gives
\[
\mathfrak H_k\to0\iff\rho_k\to0.
\]
Combining these with the numerical consistency criterion proves both consistency statements.

The same sandwiches transfer eventual comparability:
\[
\mathfrak R_k\asymp t_k^{-1}
\iff\rho_k\asymp t_k^{-1},
\qquad
\mathfrak H_k\asymp t_k^{-1/2}
\iff\rho_k\asymp t_k^{-1}.
\]
For risk, one combines either comparison with
\(\rho_k\le\mathfrak R_k/c\) and
\(\mathfrak R_k\le C_0\rho_k\).
For length, one first uses the comparison with \(\sqrt{\rho_k}\) and then squares, or takes square roots, of the nonnegative inequalities. The resource and bounded-dimension criteria therefore transfer exactly as stated.

To retain common constants for the final risk and length comparisons, define
\[
\chi_d=\frac{\log\log(ed)}{\log(ed)},\qquad
c_{\mathrm{seq}}=\min\{c/4,c^2/4\},\qquad
C_{\mathrm{seq}}=\max\{4C_0,4C_0^2\}.
\]
The numerical comparison in the preceding step and the finite bounds imply
\[
c_{\mathrm{seq}}\chi_d^2
\le\mathfrak R_{\mathcal C}(n,d,\varepsilon)
\le C_{\mathrm{seq}}\chi_d^2,
\]
and
\[
\sqrt{c_{\mathrm{seq}}}\,|\chi_d|
\le\mathfrak H_{\mathcal C}(n,d,\varepsilon)
\le\sqrt{C_{\mathrm{seq}}}\,|\chi_d|.
\]
Indeed, taking square roots of the numerical rate bounds gives
\[
\frac c2|\chi_d|
\le c\sqrt{\rho(t,d)}
\le\mathfrak H_{\mathcal C},
\qquad
\mathfrak H_{\mathcal C}
\le C_0\sqrt{\rho(t,d)}
\le2C_0|\chi_d|,
\]
and the definitions of \(c_{\mathrm{seq}},C_{\mathrm{seq}}\) encompass these factors as well as the risk factors.
For \(d\ge2\), \(\chi_d>0\).
Along a sequence with \(d_k\to\infty\) and
\(t_k/d_k^2\in[a,b_{\mathrm{ratio}}]\), the dimension threshold is eventually met.
The displayed bounds prove the two final asymptotic orders.
% lean: CausalSmith.Stat.LdpOptvalueUniformFrontier.value_sequence_conclusions_of_bounds
\end{enumerate}
\end{proof}

\begin{proof}[Proof of \cref{obj:thm:two-cell-calibration}]
\begin{enumerate}
\item
By \cref{obj:lem:measurable-kernel-radon-nikodym}, the measurable-density hypothesis of \cref{obj:lem:dimension-free-private-two-point} is satisfied. The latter supplies a universal constant \(c_1>0\) such that, for every allowed two-cell resource tuple and sampling scheme, and every \(\mathcal C\in\{\mathrm{NI},\mathrm{SI}\}\),
\[
\mathfrak R_{\mathcal C}(n,2,\varepsilon)
\ge c_1\min\left\{1,\frac1{n\varepsilon^2}\right\},
\qquad
\mathfrak H_{\mathcal C}(n,2,\varepsilon)
\ge c_1\min\left\{1,\frac1{\sqrt n\,\varepsilon}\right\}.
\]
Set
\[
c=\min\{c_1,1\},
\qquad C_0=1000.
\]
Then \(0<c\le1\le C_0<\infty\).

Fix \(n\ge2\), \(0<\varepsilon\le1\), and a sampling scheme satisfying
\cref{obj:ass:iid-people,obj:ass:independent-randomness}. Define the effective private resource by
\[
t=n\varepsilon^2>0.
\]
All expectations and probabilities below use the decision law induced by this sampling scheme and the specified protocol. We first construct simultaneous risk and interval witnesses, splitting according to \(t\le1\) or \(t>1\).
% lean: CausalSmith.Stat.LdpOptvalueUniformFrontier.two_cell_calibration

\item
Suppose \(t\le1\). Since
\[
(\sqrt n\,\varepsilon)^2=t,
\qquad 0<\sqrt n\,\varepsilon\le1,
\]
the two rate factors satisfy
\[
\min\left\{1,\frac1{n\varepsilon^2}\right\}
=
\min\left\{1,\frac1{\sqrt n\,\varepsilon}\right\}
=1.
\]
Use a singleton public seed and let every participant send the same fixed message. Set
\[
T(\mathbf Z,R,U)=\frac12,
\qquad
I(\mathbf Z,R,U)=\left[\frac14,\frac34\right].
\]
The message law is constant in the original record and history, so it satisfies full-record privacy and noninteraction. The estimate and interval endpoints are constant Borel functions.

For every \(P\in\mathcal V_2\), the interior means in
\cref{obj:def:causal-model} give
\[
\frac14\le
\max\{\mu_{0j}(P),\mu_{1j}(P)\}
\le\frac34
\qquad(j\in[2]).
\]
Averaging these inequalities yields
\[
\frac14\le
V(P)=\frac12\sum_{j=1}^{2}
\max\{\mu_{0j}(P),\mu_{1j}(P)\}
\le\frac34.
\]
Consequently,
\[
\mathbb E[(T-V(P))^2]
=\left(\frac12-V(P)\right)^2
\le\frac1{16},
\qquad
\mathbb P\{V(P)\in I\}=1,
\qquad
\mathbb E[\operatorname{length}I]=\frac12.
\]
These bounds attain both rate factors with comparison constant \(1\). Enlarging that constant to \(1000\) preserves the same witnesses and coverage.
% lean: CausalSmith.Stat.LdpOptvalueUniformFrontier.two_cell_saturated_attainment
% lean: CausalSmith.Stat.LdpOptvalueUniformFrontier.causalAttainment_mono_constant

\item
Suppose \(t>1\). We construct the three transcript statistics used in this branch. Define the binary and vector block sizes, privacy attenuation, and scaled-message magnitude by
\[
m_{\mathrm{bin}}=\left\lfloor\frac n2\right\rfloor,
\qquad
m_{\mathrm{vec}}=n-m_{\mathrm{bin}},
\qquad
\delta=\tanh(\varepsilon/2)
=\frac{e^\varepsilon-1}{e^\varepsilon+1},
\qquad
b=\frac2\delta.
\]
Here \(m_{\mathrm{bin}}\ge1\), \(m_{\mathrm{vec}}\ge1\), and \(0<\delta<1\).
For later calibration, the elementary inequalities
\[
e^\varepsilon-1\ge\varepsilon,
\qquad e^\varepsilon\le e<3
\]
also give
\[
\delta
=\frac{e^\varepsilon-1}{e^\varepsilon+1}
\ge\frac{\varepsilon}{4}.
\]

Use a singleton public seed. Each participant's message belongs to the finite alphabet
\(\{-1,1\}\times\{-1,1\}^{2}\). For an original record
\(O=(j,A,Y)\), with \(j\in[2]\) and \(A,Y\in\{0,1\}\), define, for
\(h\in\{-1,1\}\), \(z\in\{-1,1\}^{2}\), and \(1\le i\le n\),
\[
Q_i\bigl(\{(h,z)\}\mid O=(j,A,Y)\bigr)
=
\begin{cases}
\displaystyle
\frac{1+h\delta(2Y-1)}2\,
\mathbf1_{\{z=(-1,-1)\}},
& i\le m_{\mathrm{bin}},\\[6pt]
\displaystyle
\frac14\bigl(1+\delta(2A-1)(2Y-1)z_j\bigr)\,
\mathbf1_{\{h=-1\}},
& i>m_{\mathrm{bin}}.
\end{cases}
\]
These kernels are defined on every original record and are independent of message history. Their masses are nonnegative and sum to one: the first branch sums over the two binary signs, while in the second branch
\[
\sum_{z\in\{-1,1\}^{2}}z_j=0.
\]

For the binary branch, its two possible nonzero atom masses lie between
\((1-\delta)/2\) and \((1+\delta)/2\). The identity
\[
1+\delta=e^\varepsilon(1-\delta)
\]
therefore gives the atomwise privacy inequality for every pair of original records. The vector branch is precisely the \(d=2\) channel in
\cref{obj:lem:private-moment-and-dependence}, which supplies its full-record privacy. The fixed extra message coordinates preserve these inequalities. Summing over atoms gives
\[
Q_i(E\mid o)\le e^\varepsilon Q_i(E\mid o')
\]
for every message event \(E\) and every pair of original records \(o,o'\). Thus
\[
Q\in\mathfrak Q_{\mathrm{NI}}(n,2,\varepsilon).
\]

Write the transcript and its scaled vector coordinates as
\[
\mathbf Z
=\bigl((\mathsf H_i^{\mathrm{cal}},Z_i^{\mathrm{cal}})\bigr)_{i=1}^{n},
\qquad
W_{ij}^{\mathrm{cal}}
=bZ_{m_{\mathrm{bin}}+i,j}^{\mathrm{cal}}
\quad
(1\le i\le m_{\mathrm{vec}},\ j\in[2]).
\]
Define the baseline and contrast estimates by
\[
\widehat B_{\mathrm{cal}}
=\frac12+
\frac1{2\delta m_{\mathrm{bin}}}
\sum_{i=1}^{m_{\mathrm{bin}}}\mathsf H_i^{\mathrm{cal}},
\qquad
\widehat\tau_j
=\frac1{m_{\mathrm{vec}}}
\sum_{i=1}^{m_{\mathrm{vec}}}W_{ij}^{\mathrm{cal}}
\quad(j\in[2]).
\]

Conditional on the records, the transcript atom probability is the product of its participant-specific atom probabilities. Integrating against the product record law from
\cref{obj:ass:iid-people}, and using
\cref{obj:ass:independent-randomness}, gives the transcript identity
\[
\mathbb P\{\mathbf Z=((h_i,z_i))_{i=1}^{n}\}
=
\prod_{i=1}^{n}
\int Q_i(\{(h_i,z_i)\}\mid o)\,P_O(do).
\]
Thus distinct message rows are independent under the sampling decision law.

For \(i\le m_{\mathrm{bin}}\), direct summation over the binary message gives
\[
\mathbb E[\mathsf H_i^{\mathrm{cal}}\mid O_i]
=\delta(2Y_i-1),
\qquad
\mathbb E[\mathsf H_i^{\mathrm{cal}}]
=\delta(2B(P)-1),
\qquad
\mathbb E[(\mathsf H_i^{\mathrm{cal}})^2]=1.
\]
It follows that
\[
\mathbb E[\widehat B_{\mathrm{cal}}]=B(P),
\qquad
\operatorname{Var}\left(
\sum_{i=1}^{m_{\mathrm{bin}}}\mathsf H_i^{\mathrm{cal}}
\right)
=
\sum_{i=1}^{m_{\mathrm{bin}}}
\operatorname{Var}(\mathsf H_i^{\mathrm{cal}})
\le m_{\mathrm{bin}}.
\]
Scaling this centered sum proves
\[
\mathbb E[(\widehat B_{\mathrm{cal}}-B(P))^2]
\le\frac1{4m_{\mathrm{bin}}\delta^2}.
\]

For the vector block, \cref{obj:lem:private-moment-and-dependence} supplies
\[
\mathbb E[W_{ij}^{\mathrm{cal}}]=\tau_j(P),
\qquad
\mathbb E[(W_{ij}^{\mathrm{cal}})^2]=b^2.
\]
The hypotheses of that result hold with block size \(m_{\mathrm{vec}}\), since
\(0<m_{\mathrm{vec}}\le n\), \(d=2\), and \(P\in\mathcal V_2\).
Expanding the squared column sum and using independence of distinct rows yields, for each \(j\in[2]\),
\[
\begin{aligned}
\mathbb E[\widehat\tau_j]
&=\tau_j(P),\\
\mathbb E[\widehat\tau_j^2]
&=\frac{
m_{\mathrm{vec}}b^2+
m_{\mathrm{vec}}(m_{\mathrm{vec}}-1)\tau_j(P)^2
}{m_{\mathrm{vec}}^2}
=\tau_j(P)^2+
\frac{b^2-\tau_j(P)^2}{m_{\mathrm{vec}}}.
\end{aligned}
\]
Subtracting the squared mean therefore gives
\[
\mathbb E[(\widehat\tau_j-\tau_j(P))^2]
=
\frac{b^2-\tau_j(P)^2}{m_{\mathrm{vec}}}
\le\frac{b^2}{m_{\mathrm{vec}}}
=\frac4{m_{\mathrm{vec}}\delta^2}.
\]
All three statistics are Borel functions of the finite transcript alone. These are the required component bounds, simultaneously for every \(P\in\mathcal V_2\).
% lean: CausalSmith.Stat.LdpOptvalueUniformFrontier.twoCellComponents_of_sampling

\item
In the present branch,
\[
1<\sqrt n\,\varepsilon,
\qquad
\min\left\{1,\frac1{n\varepsilon^2}\right\}
=\frac1t,
\qquad
\min\left\{1,\frac1{\sqrt n\,\varepsilon}\right\}
=\frac1{\sqrt n\,\varepsilon}.
\]
Define the unprojected and projected welfare estimates by
\[
\widetilde V_{\mathrm{cal}}
=\widehat B_{\mathrm{cal}}
+\frac{|\widehat\tau_1|+|\widehat\tau_2|}{4},
\qquad
T(\mathbf Z,R,U)
=\max\left\{\frac14,
\min\left\{\frac34,\widetilde V_{\mathrm{cal}}\right\}\right\}.
\]
This \(T\) is a Borel function of \(\mathbf Z\) alone.
By \cref{obj:prop:signed-causal-bridge},
\[
V(P)=B(P)+\frac{|\tau_1(P)|+|\tau_2(P)|}{4}.
\]
The interior means also give \(V(P)\in[1/4,3/4]\), as established in the preceding branch.

Projection toward this interval reduces squared error:
\[
(T-V(P))^2
\le(\widetilde V_{\mathrm{cal}}-V(P))^2.
\]
Indeed, if \(\widetilde V_{\mathrm{cal}}\le1/4\), then
\(T=1/4\le V(P)\); if \(\widetilde V_{\mathrm{cal}}\ge3/4\), then
\(V(P)\le T=3/4\); and inside the interval \(T=\widetilde V_{\mathrm{cal}}\).

The reverse triangle inequality gives
\[
\bigl(|\widehat\tau_j|-|\tau_j(P)|\bigr)^2
\le(\widehat\tau_j-\tau_j(P))^2
\qquad(j\in[2]).
\]
Applying the squared-sum inequality first to the baseline and contrast contributions and then to the two contrast contributions proves the pointwise bound
\[
(T-V(P))^2
\le
2(\widehat B_{\mathrm{cal}}-B(P))^2
+\frac{
(\widehat\tau_1-\tau_1(P))^2+
(\widehat\tau_2-\tau_2(P))^2
}{4}.
\]
Integrating and substituting the component bounds gives
\[
\mathbb E[(T-V(P))^2]
\le
\frac1{2m_{\mathrm{bin}}\delta^2}
+\frac2{m_{\mathrm{vec}}\delta^2}.
\]

The half split satisfies
\[
m_{\mathrm{bin}}\ge\frac n3,
\qquad
m_{\mathrm{vec}}\ge\frac n2.
\]
For the first inequality, even \(n\) gives \(m_{\mathrm{bin}}=n/2\); odd \(n\ge3\) gives
\(m_{\mathrm{bin}}=(n-1)/2\ge n/3\).
The second follows from \(m_{\mathrm{bin}}\le n/2\).
Together with \(\delta\ge\varepsilon/4\), these bounds yield the numerical calibration
\[
\begin{aligned}
\mathbb E[(T-V(P))^2]
&\le
\frac1{2(n/3)\delta^2}
+\frac2{(n/2)\delta^2}\\
&=\frac{88}{n(4\delta)^2}
\le\frac{88}{n\varepsilon^2}
=\frac{88}{t}
\le\frac{1000}{t}.
\end{aligned}
\]
% lean: CausalSmith.Stat.LdpOptvalueUniformFrontier.two_cell_risk_of_component_mse

\item
Using the same protocol and estimator, define the deterministic interval radius and interval by
\[
q_{\mathrm{cal}}
=\sqrt{\frac{10000}{t}}
=\frac{100}{\sqrt n\,\varepsilon}>0,
\]
\[
I(\mathbf Z,R,U)
=
\left[
\max\left\{\frac14,T-q_{\mathrm{cal}}\right\},
\min\left\{\frac34,T+q_{\mathrm{cal}}\right\}
\right].
\]
Both endpoints are Borel functions of the transcript. Since
\(T\in[1/4,3/4]\) and \(q_{\mathrm{cal}}\ge0\),
\[
\frac14
\le\max\left\{\frac14,T-q_{\mathrm{cal}}\right\}
\le T
\le\min\left\{\frac34,T+q_{\mathrm{cal}}\right\}
\le\frac34.
\]
Thus the endpoints are ordered and the interval is connected, closed, and contained in the required domain.

For \(V(P)\in[1/4,3/4]\), the inequality
\(|T-V(P)|\le q_{\mathrm{cal}}\) implies \(V(P)\in I\). Hence
\[
\{V(P)\notin I\}
\subseteq
\{(T-V(P))^2\ge q_{\mathrm{cal}}^2\}.
\]
Markov's inequality and the risk bound give
\[
\begin{aligned}
q_{\mathrm{cal}}^2\,
\mathbb P\{(T-V(P))^2\ge q_{\mathrm{cal}}^2\}
&\le\mathbb E[(T-V(P))^2]\\
&\le\frac{88}{t}
\le\frac{1000}{t}
=\frac{q_{\mathrm{cal}}^2}{10}.
\end{aligned}
\]
Division by \(q_{\mathrm{cal}}^2>0\) and passage to the coverage complement therefore yield
\[
\mathbb P\{V(P)\in I\}\ge1-\frac1{10}=0.90.
\]

Finally, the upper endpoint is at most \(T+q_{\mathrm{cal}}\), and the lower endpoint is at least \(T-q_{\mathrm{cal}}\). Consequently,
\[
\operatorname{length}I
\le2q_{\mathrm{cal}}
=\frac{200}{\sqrt n\,\varepsilon},
\]
and, under the probability decision law,
\[
\mathbb E[\operatorname{length}I]
\le\frac{200}{\sqrt n\,\varepsilon}
\le\frac{1000}{\sqrt n\,\varepsilon}.
\]
This completes the simultaneous attainment bounds in the branch \(t>1\).
% lean: CausalSmith.Stat.LdpOptvalueUniformFrontier.twoCellInterval_coverage_of_risk
% lean: CausalSmith.Stat.LdpOptvalueUniformFrontier.twoCellInterval_expectedLength_le
% lean: CausalSmith.Stat.LdpOptvalueUniformFrontier.two_cell_unsaturated_attainment_of_components

\item
The two branches provide, for every allowed tuple and sampling scheme, a total noninteractive protocol on the original records and transcript-based decisions satisfying
\[
\sup_{P\in\mathcal V_2}\mathbb E[(T-V(P))^2]
\le1000\min\left\{1,\frac1{n\varepsilon^2}\right\},
\]
\[
\inf_{P\in\mathcal V_2}\mathbb P\{V(P)\in I\}\ge0.90,
\qquad
\sup_{P\in\mathcal V_2}\mathbb E[\operatorname{length}I]
\le1000\min\left\{1,\frac1{\sqrt n\,\varepsilon}\right\}.
\]
The definitions in
\cref{obj:def:noninteractive-class,obj:def:sequential-class} give
\[
\mathfrak Q_{\mathrm{NI}}(n,2,\varepsilon)
\subseteq
\mathfrak Q_{\mathrm{SI}}(n,2,\varepsilon).
\]
Thus these same witnesses are admissible for either protocol class. Taking the infima in
\cref{obj:def:risk,obj:def:honest-length} proves, for every
\(\mathcal C\in\{\mathrm{NI},\mathrm{SI}\}\),
\[
\mathfrak R_{\mathcal C}(n,2,\varepsilon)
\le1000\min\left\{1,\frac1{n\varepsilon^2}\right\},
\qquad
\mathfrak H_{\mathcal C}(n,2,\varepsilon)
\le1000\min\left\{1,\frac1{\sqrt n\,\varepsilon}\right\}.
\]
% lean: CausalSmith.Stat.LdpOptvalueUniformFrontier.causalAttainment_minimax_upper

\item
Apply the lower bounds obtained at the start to the fixed sampling scheme and either protocol class. Since \(c\le c_1\) and both rate factors are nonnegative,
\[
c\min\left\{1,\frac1{n\varepsilon^2}\right\}
\le
c_1\min\left\{1,\frac1{n\varepsilon^2}\right\}
\le\mathfrak R_{\mathcal C}(n,2,\varepsilon),
\]
\[
c\min\left\{1,\frac1{\sqrt n\,\varepsilon}\right\}
\le
c_1\min\left\{1,\frac1{\sqrt n\,\varepsilon}\right\}
\le\mathfrak H_{\mathcal C}(n,2,\varepsilon).
\]
Together with the upper bounds and the simultaneous witnesses, these prove all assertions with the chosen universal constants \(c\) and \(C_0\).
% lean: CausalSmith.Stat.LdpOptvalueUniformFrontier.two_cell_calibration
\end{enumerate}
\end{proof}

\begin{proof}[Proof of \cref{obj:thm:paired-uniform-tv-frontier}]
Throughout, resources are admissible:
\[
n,d\in\mathbb N,\qquad n,d\ge2,\qquad 0<\varepsilon\le1.
\]
We use the effective resource, logarithmic dimension, and converse tuning quantities
\[
t=n\varepsilon^2>0,\qquad L=\log(ed),\qquad
z_*=\frac{d^2L}{t},\qquad w=\log(e+z_*).
\]
In particular, \(1\le L\le d\). The rates in \cref{obj:def:frontier-handle} are independent of the protocol-class label:
\[
r_{\mathrm{NI}}=r_{\mathrm{SI}}=\rho(t,d),
\qquad
h_{\mathrm{NI}}=h_{\mathrm{SI}}=\sqrt{\rho(t,d)}.
\]

\begin{enumerate}
\item \textit{A uniform approximation constant.}
Invoke the published Bernstein limit
\[
K e_K\longrightarrow\beta_*
\quad\text{along positive even integers }K,
\qquad 0.278<\beta_*<0.286,
\]
as stated in \citep[Section 3.1, definition of the Bernstein constant and the immediately following paragraph, p.~8]{CaiLow2011Nonsmooth}; see also \citep{Bernstein1913}.
Choose an integer \(N_{\mathrm B}\ge1\) such that
\[
K e_K\ge\frac1{10}
\qquad\text{for every even }K\ge2N_{\mathrm B},
\]
and define
\[
b_{\mathrm{app}}
=\min\left\{\frac1{10},e_{2N_{\mathrm B}}\right\}>0.
\]
The positivity follows from the preceding tail bound at \(K=2N_{\mathrm B}\).
Approximation error decreases as the allowed degree increases. Hence, for every positive even \(K\),
\[
\frac{b_{\mathrm{app}}}{K}\le e_K.
\]
Indeed, the tail bound proves this when \(K\ge2N_{\mathrm B}\); otherwise,
\[
b_{\mathrm{app}}\le e_{2N_{\mathrm B}}\le e_K\le K e_K.
\]
% lean: CausalSmith.Stat.LdpOptvalueUniformFrontier.bernstein_uniform_positive_of_gate

\item \textit{Paired product priors.}
Take the converse amplitude and matching degree from \cref{obj:def:frontier-handle}:
\[
(\alpha,k)=
\begin{cases}
\left(\dfrac{\sqrt{z_*}}{100},\,2\lceil16L\rceil+1\right),&z_*\le1,\\[4pt]
\left(\dfrac1{100},\,2\lceil16L/w\rceil+1\right),&z_*>1,
\end{cases}
\qquad K=k-1.
\]
These choices satisfy \(0<\alpha\le1/2\), and \(K\) is a positive even integer.
By \cref{obj:lem:cai-low-moment-duality}, there are symmetric probability measures \(\xi_0,\xi_1\), supported on \([-1,1]\), whose moments agree through degree \(K\) and whose absolute moments differ by \(2e_K\). Define the scaled coordinate priors and their product priors by
\[
\nu_\ell=(x\mapsto\alpha x)_\#\xi_\ell,
\qquad
\Pi_\ell=\nu_\ell^{\otimes d},
\qquad \ell\in\{0,1\}.
\]
They are supported on \([-\alpha,\alpha]^d\subseteq[-1/2,1/2]^d\), and
\[
\int u^q\,d\nu_1(u)=\int u^q\,d\nu_0(u)
\qquad(0\le q\le K).
\]
Using the paired-target identity in \cref{obj:prop:signed-causal-bridge}, define their total-variation centers and separation by
\[
v_\ell^{\mathrm{TV}}=\frac12\int |u|\,d\nu_\ell(u),
\qquad
\Delta_{\mathrm{TV}}
=v_1^{\mathrm{TV}}-v_0^{\mathrm{TV}}
=\alpha e_K>0.
\]
Independence of the coordinates gives
\[
\mathbb E_{\Pi_\ell}D_{\mathrm{TV}}(P_\theta)=v_\ell^{\mathrm{TV}},
\qquad
\operatorname{Var}_{\Pi_\ell}\!\left(D_{\mathrm{TV}}(P_\theta)\right)
\le\frac{\alpha^2}{4d}.
\]
Consequently, Chebyshev's inequality and the approximation bound give
\[
\Pi_\ell\left\{
\left|D_{\mathrm{TV}}(P_\theta)-v_\ell^{\mathrm{TV}}\right|
\ge\frac{\Delta_{\mathrm{TV}}}{8}
\right\}
\le
\frac{16\alpha^2}{d\Delta_{\mathrm{TV}}^2}
\le\frac{16K^2}{b_{\mathrm{app}}^2d}.
\]

The ceiling inequalities also give the required rate calibration. If \(z_*\le1\), then
\[
K\le34L,\qquad
\sqrt{\rho(t,d)}=\frac{\sqrt{z_*}}{L},
\qquad
\frac{\alpha}{K}\ge\frac1{3400}\sqrt{\rho(t,d)}.
\]
If \(z_*>1\), then
\[
K\le34\max\{1,L/w\},
\qquad
\sqrt{\rho(t,d)}=\min\{1,w/L\},
\]
so again
\[
\frac{\alpha}{K}\ge\frac1{3400}\sqrt{\rho(t,d)}.
\]
Since \(w\ge1\), both branches satisfy \(K\le34L\). Define
\[
s_{\mathrm{app}}=\frac{b_{\mathrm{app}}}{3400}>0.
\]
Then
\[
\Delta_{\mathrm{TV}}\ge s_{\mathrm{app}}\sqrt{\rho(t,d)}.
\]
Finally, \(L^2/d\to0\) as \(d\to\infty\). Thus there is a universal integer \(D_{\mathrm{cut}}\ge2\) such that
\[
\frac{16\cdot34^2L^2}{b_{\mathrm{app}}^2d}\le\frac1{100}
\qquad(d>D_{\mathrm{cut}}).
\]
For those dimensions, each prior's target-concentration failure probability is at most \(1/100\).
% lean: CausalSmith.Stat.LdpOptvalueUniformFrontier.paired_lawspace_resource_priors_of_gate

\item \textit{Contraction of the paired transcript mixtures.}
Fix any sequential paired-input protocol \(\mathsf K\).
By \cref{obj:prop:signed-causal-bridge}, its pullback
\[
Q_i(E\mid(j,a,y),\eta,r)
=
\mathsf K_i\!\left(E\mid(j,(2a-1)(2y-1)),\eta,r\right)
\]
is sequentially private at the same budget, and its seed-transcript law under \(G_\theta\) equals that of \(\mathsf K\) under \(P_\theta\).
The density hypothesis for \cref{obj:lem:adaptive-moment-contraction} is supplied by \cref{obj:lem:measurable-kernel-radon-nikodym}.
Define the paired seed-transcript mixtures by
\[
\mathsf M_\ell^{\mathrm{TV}}
=
\int\mathcal L_{P_\theta,\mathsf K}(R,\mathbf Z^{\mathsf K})\,
d\Pi_\ell(\theta),
\qquad \ell\in\{0,1\},
\]
and the contraction scale by
\[
\beta=\frac{e^\varepsilon-1}{2d}.
\]
The moment agreement in the previous step allows \cref{obj:lem:adaptive-moment-contraction} to yield
\[
\operatorname{TV}(\mathsf M_0^{\mathrm{TV}},\mathsf M_1^{\mathrm{TV}})
\le
\min\left\{1,d\alpha^k\sqrt{\binom nk}\,\beta^k\right\}
\qquad(k\le n),
\]
whereas the two mixtures are equal when \(k>n\).

For \(0<\varepsilon\le1\), the exponential bound
\(e^\varepsilon-1\le2\varepsilon\) gives \(\beta\le\varepsilon/d\).
Moreover,
\[
\binom nk\le\frac{n^k}{k!}\le\left(\frac{en}{k}\right)^k,
\]
because \(k^k/k!\le e^k\), as follows from the exponential series.
Thus, when \(k\le n\),
\[
\operatorname{TV}(\mathsf M_0^{\mathrm{TV}},\mathsf M_1^{\mathrm{TV}})
\le
\min\left\{1,
d\left(\alpha\sqrt{\frac{en}{k}}\frac{\varepsilon}{d}\right)^k
\right\}.
\]

To evaluate this bound, define
\[
w_{\mathrm{con}}=
\begin{cases}
1,&z_*\le1,\\
w,&z_*>1,
\end{cases}
\qquad
B_{\mathrm{con}}=\alpha\sqrt{\frac{en}{k}}\frac{\varepsilon}{d}.
\]
The prescribed ceilings give
\[
32L\le k w_{\mathrm{con}},
\qquad
B_{\mathrm{con}}^2=\frac{\alpha^2eL}{kz_*}.
\]
In the branch \(z_*\le1\), the amplitude choice and \(e<3\) imply
\[
\alpha^2e\,w_{\mathrm{con}}e^{w_{\mathrm{con}}/2}
=\frac{z_*}{10^4}e^{3/2}\le32z_*.
\]
In the other branch, \(e^w=e+z_*\), and \(w\le2e^{w/2}\) gives
\[
w e^{w/2}\le2(e+z_*)\le8z_*.
\]
Hence the same bound
\[
\alpha^2e\,w_{\mathrm{con}}e^{w_{\mathrm{con}}/2}\le32z_*
\]
holds. Combining it with \(32L\le kw_{\mathrm{con}}\) yields
\[
B_{\mathrm{con}}^2\le e^{-w_{\mathrm{con}}/2},
\qquad
B_{\mathrm{con}}\le e^{-w_{\mathrm{con}}/4}.
\]
Therefore
\[
d B_{\mathrm{con}}^k
\le d e^{-8L}
\le e^{-7L}
\le e^{-7}\le\frac1{100}.
\]
For \(k>n\), exact equality of the mixtures gives the same conclusion:
\[
\operatorname{TV}(\mathsf M_0^{\mathrm{TV}},\mathsf M_1^{\mathrm{TV}})
\le\frac1{100}.
\]
% lean: CausalSmith.Stat.LdpOptvalueUniformFrontier.paired_lawspace_resource_contraction

\item \textit{Decision lower bounds in large dimension.}
We give the paired-model decision reduction explicitly.
For an estimator and a globally \(0.90\)-honest interval, define their worst-law risk and expected length by
\[
\begin{aligned}
\mathcal R_{\mathsf K}(\widehat D)
&=\sup_{\mathsf p\in\mathcal P_d^{\mathrm{pair}}}
\mathbb E_{\mathsf p,\mathsf K}
\left[(\widehat D-D_{\mathrm{TV}}(\mathsf p))^2\right],\\
\mathcal H_{\mathsf K}(J_{\mathrm{TV}})
&=\sup_{\mathsf p\in\mathcal P_d^{\mathrm{pair}}}
\mathbb E_{\mathsf p,\mathsf K}\operatorname{length}J_{\mathrm{TV}}.
\end{aligned}
\]
These definitions apply to every admissible Borel decision of \cref{obj:def:tv-private-decision-values}.
For this reduction, let \(\eta_0,\omega\ge0\) bound the two concentration failures and the seed-transcript total variation, respectively.
Let
\[
N_\ell^{\mathrm{TV}}
=\mathsf M_\ell^{\mathrm{TV}}\otimes\operatorname{Uniform}[0,1]
\]
be the joint seed-transcript-randomizer mixture.
Appending the independent analyst randomizer preserves the bound:
\[
\operatorname{TV}(N_0^{\mathrm{TV}},N_1^{\mathrm{TV}})\le\omega.
\]
Indeed, integrating an event indicator over the independent randomizer produces a measurable function taking values in \([0,1]\).

Threshold \(\widehat D\) at
\((v_0^{\mathrm{TV}}+v_1^{\mathrm{TV}})/2\).
On the target-concentration event, a wrong prior label entails an estimation error of magnitude at least \(3\Delta_{\mathrm{TV}}/8\).
If the worst-law risk is finite, Markov's inequality bounds the sum of testing errors by
\[
2\eta_0+
\frac{128\,\mathcal R_{\mathsf K}(\widehat D)}
{9\Delta_{\mathrm{TV}}^2}.
\]
The testing lower bound from total variation is \(1-\omega\). It follows that
\[
\mathcal R_{\mathsf K}(\widehat D)
\ge
\frac{9\Delta_{\mathrm{TV}}^2}{128}
(1-\omega-2\eta_0)_+.
\]
The same inequality holds immediately when the risk is infinite.

For the interval reduction, define
\[
\mathcal E_\ell^{\mathrm{TV}}
=
\left\{
J_{\mathrm{TV}}\cap
\left[v_\ell^{\mathrm{TV}}-\frac{\Delta_{\mathrm{TV}}}{8},
      v_\ell^{\mathrm{TV}}+\frac{\Delta_{\mathrm{TV}}}{8}\right]
\ne\varnothing
\right\}.
\]
Paired-model honesty and target concentration imply
\[
N_\ell^{\mathrm{TV}}(\mathcal E_\ell^{\mathrm{TV}})
\ge0.90-\eta_0.
\]
Transferring the second event to the first mixture and applying the intersection bound gives
\[
N_0^{\mathrm{TV}}
(\mathcal E_0^{\mathrm{TV}}\cap\mathcal E_1^{\mathrm{TV}})
\ge0.80-2\eta_0-\omega.
\]
On this intersection, ordered endpoints and connectedness force length at least \(3\Delta_{\mathrm{TV}}/4\). Therefore
\[
\mathcal H_{\mathsf K}(J_{\mathrm{TV}})
\ge
\frac{3\Delta_{\mathrm{TV}}}{4}
(0.80-2\eta_0-\omega)_+.
\]

For \(d>D_{\mathrm{cut}}\), the preceding steps permit
\(\eta_0=\omega=1/100\). Define
\[
c_{\mathrm{large}}
=\min\left\{\frac{s_{\mathrm{app}}^2}{100},
            \frac{s_{\mathrm{app}}}{100}\right\}>0.
\]
Since \(\Delta_{\mathrm{TV}}\ge s_{\mathrm{app}}\sqrt{\rho(t,d)}\),
\[
\begin{aligned}
\mathcal R_{\mathsf K}(\widehat D)
&\ge\frac9{128}\frac{97}{100}\Delta_{\mathrm{TV}}^2
\ge c_{\mathrm{large}}\rho(t,d),\\
\mathcal H_{\mathsf K}(J_{\mathrm{TV}})
&\ge\frac34\frac{77}{100}\Delta_{\mathrm{TV}}
\ge c_{\mathrm{large}}\sqrt{\rho(t,d)}.
\end{aligned}
\]
% lean: CausalSmith.Stat.LdpOptvalueUniformFrontier.paired_converse_of_gate

\item \textit{Elementary rate comparisons and bounded dimension.}
We first establish the numerical comparisons used to complete the converse:
\[
\min\{1,t^{-1}\}\le\rho(t,d)\le4\min\{1,d^2/t\},
\qquad
\frac14\min\{t,d\}\le t\rho(t,d).
\]
In the dense branch \(t\ge d^2L\), these follow from
\[
\rho(t,d)=\frac{d^2}{tL},\qquad
\frac1t\le\rho(t,d)\le\min\{1,d^2/t\},
\qquad
t\rho(t,d)=\frac{d^2}{L}\ge d.
\]

For the nondense branch, define the auxiliary function
\[
\Psi(s_{\mathrm{res}})
=s_{\mathrm{res}}
\left[\log\left(e+\frac{d^2L}{s_{\mathrm{res}}}\right)\right]^2,
\qquad s_{\mathrm{res}}>0.
\]
Writing \(v_{\mathrm{res}}=d^2L/s_{\mathrm{res}}>0\), differentiation gives
\[
\Psi'(s_{\mathrm{res}})
=
\log(e+v_{\mathrm{res}})
\left[
\log(e+v_{\mathrm{res}})
-\frac{2v_{\mathrm{res}}}{e+v_{\mathrm{res}}}
\right]\ge0.
\]
For the last inequality, \(\log y\ge1-y^{-1}\), applied at
\(y=(e+v_{\mathrm{res}})/e\), gives
\[
\log(e+v_{\mathrm{res}})
\ge1+\frac{v_{\mathrm{res}}}{e+v_{\mathrm{res}}}
\ge\frac{2v_{\mathrm{res}}}{e+v_{\mathrm{res}}}.
\]
Also,
\[
\log(e+d^2L)\ge L,
\qquad
\log(e+dL)\ge L/2.
\]
The first follows from \(dL\ge e\); the second follows by combining
\(\log(e+dL)\ge1\) and \(\log(e+dL)\ge\log d\).
Thus, if \(t\le1\), then \(w\ge L\) and \(\rho(t,d)=1\).
If \(t\ge1\), monotonicity of \(\Psi\) gives
\(tw^2\ge L^2\), establishing the lower comparison with \(\min\{1,t^{-1}\}\).
Likewise, if \(t\le d\), then \(w\ge L/2\); if \(t\ge d\), monotonicity gives
\(tw^2\ge dL^2/4\). Applying the cap at one proves the comparison with \(\min\{t,d\}/4\).

For the upper comparison, \(z_*>1\) in this branch and
\[
w\le\log(e+1)+\log z_*
\le2+2(\sqrt{z_*}-1)=2\sqrt{z_*}.
\]
Hence
\[
\frac{w^2}{L^2}\le\frac{4z_*}{L^2}
=\frac{4d^2}{tL}\le\frac{4d^2}{t}.
\]
Together with \(\rho(t,d)\le1\), this proves all three comparisons.

Now define the two-point amplitude
\[
\alpha_{\mathrm{two}}
=\frac18\min\{1,t^{-1/2}\}\in(0,1/2].
\]
Consider the paired laws with constant contrast zero and constant contrast
\(\alpha_{\mathrm{two}}\). Their targets are zero and
\(\alpha_{\mathrm{two}}/2\).
Applying \cref{obj:lem:adaptive-moment-contraction} at matching order one through the bridge in \cref{obj:prop:signed-causal-bridge} gives
\[
\operatorname{TV}\!\left(
\mathcal L_{P_0,\mathsf K}(R,\mathbf Z^{\mathsf K}),
\mathcal L_{P_{\alpha_{\mathrm{two}}\mathbf1},\mathsf K}
(R,\mathbf Z^{\mathsf K})
\right)
\le
\alpha_{\mathrm{two}}\sqrt n\,\varepsilon
\le\frac18.
\]
Here \(P_{\alpha_{\mathrm{two}}\mathbf1}\) denotes the law with
\(\theta_j=\alpha_{\mathrm{two}}\) for every \(j\).
The preceding decision reduction, with deterministic targets and
\(\omega=1/8\), therefore yields
\[
\mathcal R_{\mathsf K}(\widehat D)
\ge\frac{63\alpha_{\mathrm{two}}^2}{4096}
\ge\frac1{8192}\min\{1,t^{-1}\},
\]
\[
\mathcal H_{\mathsf K}(J_{\mathrm{TV}})
\ge\frac{81\alpha_{\mathrm{two}}}{320}
\ge\frac1{8192}\min\{1,t^{-1/2}\}.
\]
The calibrations use
\[
\alpha_{\mathrm{two}}^2
=\frac1{64}\min\{1,t^{-1}\}.
\]

If \(d\le D_{\mathrm{cut}}\), the elementary upper comparison implies
\[
\rho(t,d)
\le4D_{\mathrm{cut}}^2\min\{1,t^{-1}\},
\]
and taking square roots gives
\[
\sqrt{\rho(t,d)}
\le2D_{\mathrm{cut}}\min\{1,t^{-1/2}\}
\le4D_{\mathrm{cut}}^2\min\{1,t^{-1/2}\}.
\]
Define
\[
c_{\mathrm{small}}=\frac1{32768D_{\mathrm{cut}}^2},
\qquad
c_{\mathrm{lower}}=\min\{c_{\mathrm{large}},c_{\mathrm{small}}\}>0.
\]
The cases \(d\le D_{\mathrm{cut}}\) and \(d>D_{\mathrm{cut}}\) now give, for every sequential paired protocol and every admissible decision,
\[
\mathcal R_{\mathsf K}(\widehat D)\ge
c_{\mathrm{lower}}\rho(t,d),
\qquad
\mathcal H_{\mathsf K}(J_{\mathrm{TV}})\ge
c_{\mathrm{lower}}\sqrt{\rho(t,d)}.
\]
Every noninteractive protocol is sequential. Taking the decision and protocol infima in \cref{obj:def:tv-private-decision-values} therefore proves these lower bounds for both classes.
% lean: CausalSmith.Stat.LdpOptvalueUniformFrontier.paired_minimax_lower_of_gate

\item \textit{The concrete upper-risk construction.}
Choose the universal constants
\[
c=\min\{c_{\mathrm{lower}},1\},
\qquad C_0=10^{16}.
\]
Then \(0<c\le C_0\).
Use the paired procedure specified in \cref{obj:synth_5}, with a singleton public seed. Its channel quantities are
\[
m=\lfloor n/3\rfloor,\qquad
\delta=\tanh(\varepsilon/2),\qquad b=\frac d\delta.
\]
Each of the first \(2m\) participants releases a vector with probabilities
\[
\mathsf K_i(z'\mid(j,s))
=2^{-d}(1+\delta s z'_j),
\qquad z'\in\{-1,1\}^d;
\]
the remaining participants release \((-1,\ldots,-1)\).
Symmetry of the sign-vector alphabet makes the active probabilities sum to one, and
\[
\frac{1+\delta}{1-\delta}=e^\varepsilon
\]
bounds the likelihood ratio at every active message. The constant releases satisfy the privacy inequality as well. All kernels are independent of history, so this is a pure \(\varepsilon\)-private noninteractive protocol with the declared sign-vector message spaces.

For \(n>2\), define the pilot and evaluation matrices, means, and calibration quantities by
\[
W^{\mathrm{pil}}_{ij}=bZ_{i,j},\qquad
W^{\mathrm{ev}}_{ij}=bZ_{m+i,j},
\qquad 1\le i\le m,\quad j\in[d],
\]
\[
M_j=\frac1m\sum_iW^{\mathrm{pil}}_{ij},\qquad
\overline W_j=\frac1m\sum_iW^{\mathrm{ev}}_{ij},
\qquad
\sigma^2=\frac{b^2}{m},\qquad \sigma=\sqrt{\sigma^2},
\qquad H=\log(e+\sigma^2L).
\]
The statistic \(\widehat F\) is the ordered branch-selected norm estimate in
\cref{obj:synth_5}, using the private moments and polynomial coefficients specified there and in \cref{obj:def:private-moments,obj:def:frontier-handle}.
Set
\[
\widehat D=
\begin{cases}
1/8,&n=2,\\
\max\{0,\min\{1/4,\widehat F/2\}\},&n>2.
\end{cases}
\]
The decisions are measurable functions of the finite transcript.

The cube target belongs to \([0,1/4]\), so
\[
(\widehat D-D_{\mathrm{TV}}(P_\theta))^2\le\frac1{16}.
\]
For \(n>2\), clipping also gives the sharper comparison
\[
\mathbb E_{P_\theta,\mathsf K}
(\widehat D-D_{\mathrm{TV}}(P_\theta))^2
\le\frac14
\mathbb E_{P_\theta,\mathsf K}(\widehat F-F(\theta))^2.
\]
To apply the existing moment and polynomial calibrations, use \(G_\theta\) from
\cref{obj:def:symmetric-law}. Its contrasts satisfy
\[
\tau_j(G_\theta)
=\frac{1+\theta_j}{2}-\frac{1-\theta_j}{2}
=\theta_j.
\]
The bridge in \cref{obj:prop:signed-causal-bridge} identifies the active paired blocks with independent vector-channel blocks under this law. Thus
\cref{obj:lem:private-moment-and-dependence,obj:lem:finite-private-polynomial-calibration}
apply with contrast vector \(\theta\); their concentration and approximation hypotheses are supplied by
\cref{obj:lem:bounded-mean-concentration,obj:lem:cai-low-chebyshev-approximation}.

We verify the resource calibration used in the risk bounds. For \(n\ge3\),
\[
m\ge n/5,\qquad
\delta=\frac{e^\varepsilon-1}{e^\varepsilon+1}\ge\varepsilon/4,
\qquad
\sigma^2\le\frac{80d^2}{t}.
\]
The attenuation bound follows from \(e^\varepsilon\ge1+\varepsilon\) and
\(e^\varepsilon+1<4\).
Consequently,
\[
1\le H\le\log(80(e+z_*))
\le5+w\le6w.
\]

We now follow the risk case split.

If \(n=2\), then \(t\le2\), and the elementary lower comparison gives
\(\rho(t,d)\ge1/2\). The constant estimate satisfies
\[
\mathbb E(\widehat D-D_{\mathrm{TV}}(P_\theta))^2
\le\frac1{64}\le C_0\rho(t,d).
\]

Suppose \(n>2\). If \(\rho(t,d)=1\), the uniform \(1/16\) loss bound gives the desired result.
Next consider the constant-fallback branch. Its conditions include
\[
t<d^2L,\qquad \frac{L}{1024H}<4.
\]
Since \(H\le6w\),
\[
L<4096H\le24576w,
\qquad
\rho(t,d)\ge\left(\frac1{24576}\right)^2.
\]
Again,
\[
\mathbb E(\widehat D-D_{\mathrm{TV}}(P_\theta))^2
\le\frac1{16}\le C_0\rho(t,d).
\]

For the remaining branches, suppose first that \(L<4096\).
In the dense case \(t\ge d^2L\), the absolute-mean statistic satisfies
\[
\mathbb E\left[
\left(\frac1d\sum_j|\overline W_j|-F(\theta)\right)^2
\right]
\le\frac1d\sum_j\mathbb E(\overline W_j-\theta_j)^2
\le\sigma^2.
\]
Here the first inequality uses the absolute-value Lipschitz inequality and the quadratic mean inequality; the second uses the moment bounds and independence across participants. The numerical calibration is
\[
\frac{24}{t}+\frac{\sigma^2}{2}
\le\frac{64d^2}{t}
=64L\rho(t,d)
\le64\cdot4096\,\rho(t,d).
\]
Since the TV risk is at most \(\sigma^2/4\), this proves its \(C_0\rho(t,d)\) bound.
In the nondense case, \(w\ge1\) implies
\[
\rho(t,d)\ge\left(\frac1{4096}\right)^2,
\]
so the uniform loss bound suffices.

Finally, suppose \(L\ge4096\).
In the dense case, the prescribed hybrid degree and radius are
\[
D=2\lfloor L/2048\rfloor,\qquad
T_*=4\sigma\sqrt L,\qquad a_{\mathrm{rad}}=8\sigma\sqrt L.
\]
Its degree fits the block:
\[
2D\le L/512\le m.
\]
Indeed, \(n\ge t\ge d^2L\ge4L\) and \(m\ge n/5\).
The hybrid conclusion of \cref{obj:lem:finite-private-polynomial-calibration} gives
\[
\mathbb E(\widehat F-F(\theta))^2
\le500000000\,\frac{\sigma^2}{L}
\le40000000000\,\rho(t,d).
\]
Also \(t^{-1}\le\rho(t,d)\), and hence
\[
\frac{24}{t}
+\frac12\left(500000000\,\frac{\sigma^2}{L}\right)
\le(24+20000000000)\rho(t,d)
\le C_0\rho(t,d).
\]
The clipping comparison proves the claimed TV risk bound.

In the nondense case remaining after the saturation and fallback cases,
\[
\rho(t,d)=(w/L)^2,\qquad
a_{\mathrm{rad}}=\frac12,\qquad
D=2\left\lfloor\frac{L}{2048H}\right\rfloor,
\]
and floor rounding gives
\[
2\le D,\qquad
\frac{L}{2048H}\le D\le\frac{L}{1024H}.
\]
The moment-fit condition \(2D\le m\) also holds. To see this, if \(m<2D\), the degree budget and \(H\ge1\) imply \(m<L/512\).
Since \(\delta<1\),
\[
\sigma^2\ge\frac{d^2}{m},
\qquad
e+\sigma^2L>512d^2\ge ed,
\]
so \(H\ge L\), contradicting \(1024DH\le L\) and \(D\ge2\).

The global-polynomial conclusion of
\cref{obj:lem:finite-private-polynomial-calibration} now yields
\[
\mathbb E(\widehat F-F(\theta))^2
\le\frac{e^{9DH}}{2d}+\frac1{(D+1)^2}.
\]
The degree budget gives \(9DH\le L/100\).
For \(L\ge4096\), the elementary exponential inequalities
\[
L^2\le16e^{L/2},
\qquad
4\le e^{49L/100-1}
\]
imply
\[
L^2e^{L/100}\le4e^{L-1}=4d.
\]
Therefore
\[
\frac{e^{9DH}}{2d}\le\frac2{L^2}\le2\rho(t,d),
\qquad
\frac1{D+1}\le\frac{2048H}{L}
\le\frac{12288w}{L},
\]
and hence
\[
\frac1{(D+1)^2}\le150994944\,\rho(t,d).
\]
The elementary lower comparison, together with \(\rho(t,d)<1\), gives
\(t^{-1}\le\rho(t,d)\). Thus
\[
\frac{24}{t}
+\frac12\left(\frac{e^{9DH}}{2d}+\frac1{(D+1)^2}\right)
\le75497500\,\rho(t,d)
\le C_0\rho(t,d).
\]
Clipping again completes the risk bound. Every admissible tuple has therefore been covered:
\[
\sup_{\mathsf p\in\mathcal P_d^{\mathrm{pair}}}
\mathbb E_{\mathsf p,\mathsf K}
(\widehat D-D_{\mathrm{TV}}(\mathsf p))^2
\le C_0\rho(t,d).
\]
% lean: CausalSmith.Stat.LdpOptvalueUniformFrontier.tvFrontierEstimator_uniform_rate

\item \textit{Intervals and the common minimax bounds.}
Define
\[
q_*=\sqrt{10C_0\rho(t,d)}=\sqrt{10^{17}\rho(t,d)},
\qquad
J_{\mathrm{TV}}
=
\left[
\max\{0,\widehat D-q_*\},
\min\{1/4,\widehat D+q_*\}
\right].
\]
The elementary lower comparison gives \(\rho(t,d)>0\).
The interval contains its center, has ordered Borel endpoints, and is connected, closed, and contained in \([0,1/4]\).
Since the target is also in this range, failure of coverage entails
\(|\widehat D-D_{\mathrm{TV}}(\mathsf p)|>q_*\). Markov's inequality gives
\[
\Pr_{\mathsf p,\mathsf K}
\{D_{\mathrm{TV}}(\mathsf p)\notin J_{\mathrm{TV}}\}
\le\frac{C_0\rho(t,d)}{q_*^2}
=\frac1{10}.
\]
Furthermore,
\[
\operatorname{length}J_{\mathrm{TV}}
\le2q_*
=2\sqrt{10C_0}\sqrt{\rho(t,d)}
\le C_0\sqrt{\rho(t,d)}.
\]
This proves all three concrete attainment guarantees.
The same construction supplies the general attainment witness: its message spaces are exactly \(\{-1,1\}^d\), and its estimate depends on the transcript alone.

The construction belongs to both protocol classes. Its interval, being contained in \([0,1/4]\), is also an admissible interval in the optimization over \([0,1]\).
Taking the corresponding candidates in the minimax infima gives, for either \(\mathcal C\),
\[
\mathfrak R^{\mathrm{TV}}_{\mathcal C}
(\mathcal P_d^{\mathrm{pair}};n,d,\varepsilon)
\le C_0\rho(t,d),
\qquad
\mathfrak H^{\mathrm{TV}}_{\mathcal C}
(\mathcal P_d^{\mathrm{pair}};n,d,\varepsilon)
\le C_0\sqrt{\rho(t,d)}.
\]
Since \(c\le c_{\mathrm{lower}}\) and both rates are nonnegative, the lower bounds proved above give
\[
\begin{aligned}
c\rho(t,d)
&\le\mathfrak R^{\mathrm{TV}}_{\mathcal C}
(\mathcal P_d^{\mathrm{pair}};n,d,\varepsilon)
\le C_0\rho(t,d),\\
c\sqrt{\rho(t,d)}
&\le\mathfrak H^{\mathrm{TV}}_{\mathcal C}
(\mathcal P_d^{\mathrm{pair}};n,d,\varepsilon)
\le C_0\sqrt{\rho(t,d)}.
\end{aligned}
\]
% lean: CausalSmith.Stat.LdpOptvalueUniformFrontier.tvAttainment_minimax_upper

\item \textit{Identification of the paired functional.}
For \(d\ge2\) and \(\theta\in[-1/2,1/2]^d\), the masses
\[
P_\theta(j,s)=\frac{1+s\theta_j}{2d}
\]
are nonnegative, and their sum is one because each pair has total mass \(1/d\).
Taking zero contrast gives the uniform probability law.
For two probability laws on a finite alphabet, total variation equals half the sum of the absolute mass differences. Therefore
\[
D_{\mathrm{TV}}(P_\theta)
=\operatorname{TV}(P_\theta,\mathsf U_{2d}).
\]
For each coordinate,
\[
\sum_{s\in\{-1,1\}}
\left|\frac{1+s\theta_j}{2d}-\frac1{2d}\right|
=
\frac{|\theta_j|}{d}.
\]
Summing and multiplying by one half gives
\[
D_{\mathrm{TV}}(P_\theta)
=\frac1{2d}\sum_{j=1}^d|\theta_j|
=\frac12F(\theta),
\qquad
0\le D_{\mathrm{TV}}(P_\theta)\le\frac14.
\]
% lean: CausalSmith.Stat.LdpOptvalueUniformFrontier.tvFromUniform_pairedLaw

\item \textit{Attainment and protocol-class ratios.}
The witnesses constructed above establish the existential attainment assertions for either class, because the class-selected rates agree. They also establish the concrete guarantees for the procedure of
\cref{obj:synth_5,obj:def:frontier-handle,obj:thm:uniform-private-value-frontiers}.

The protocol inclusion in \cref{obj:def:tv-private-decision-values} gives
\[
\mathfrak R^{\mathrm{TV}}_{\mathrm{SI}}
(\mathcal P_d^{\mathrm{pair}};n,d,\varepsilon)
\le
\mathfrak R^{\mathrm{TV}}_{\mathrm{NI}}
(\mathcal P_d^{\mathrm{pair}};n,d,\varepsilon),
\]
and the corresponding inequality for honest length.
The common bounds make all four values finite and strictly positive.
For risk, they also give
\[
\mathfrak R^{\mathrm{TV}}_{\mathrm{NI}}
(\mathcal P_d^{\mathrm{pair}};n,d,\varepsilon)
\le C_0\rho(t,d)
=\frac{C_0}{c}\,c\rho(t,d)
\le\frac{C_0}{c}
\mathfrak R^{\mathrm{TV}}_{\mathrm{SI}}
(\mathcal P_d^{\mathrm{pair}};n,d,\varepsilon).
\]
Dividing by the positive SI value yields the asserted risk ratio bounds.
Replacing \(\rho(t,d)\) by \(\sqrt{\rho(t,d)}\) proves the same ratio bounds for honest length.
% lean: CausalSmith.Stat.LdpOptvalueUniformFrontier.frontier_ratio_comparison

\item \textit{Saturation and consistency of the rate.}
For \(0<t<d^2L\), positivity of \(w/L\) gives
\[
\begin{aligned}
\rho(t,d)=1
&\Longleftrightarrow w\ge L\\
&\Longleftrightarrow e+\frac{d^2L}{t}\ge ed\\
&\Longleftrightarrow
t\le\frac{d^2L}{ed-e}.
\end{aligned}
\]
The last denominator is positive for every \(d\ge2\).

For an admissible resource sequence, define
\[
t_k=n_k\varepsilon_k^2,\qquad
L_k=\log(ed_k),\qquad
g_k=\frac{\log(1+d_k^2/t_k)}{L_k}.
\]
If \(\rho(t_k,d_k)\to0\), the lower comparison
\(\min\{1,t_k^{-1}\}\le\rho(t_k,d_k)\) forces \(t_k\to\infty\).
Also, whenever \(\rho(t_k,d_k)<1\),
\[
g_k^2\le\rho(t_k,d_k).
\]
In the nondense branch this follows from
\[
\log(1+d_k^2/t_k)
\le\log(e+L_kd_k^2/t_k)
\]
and removal of the cap when the rate is below one.
In the dense branch, using \(\log(1+x)\le x\) and
\(d_k^2/t_k\le1/L_k\) gives the same inequality.
Thus \(g_k\to0\).

Conversely, for every \(d\ge2\) and \(t>0\),
\[
\rho(t,d)\le
\left[
\frac{\log(1+d^2/t)}{L}
+\frac{\log(e+L)}{L}
\right]^2+\frac1{L^2}.
\]
For the nondense branch, this follows from
\[
e+L(d^2/t)\le(1+d^2/t)(e+L);
\]
for the dense branch, \(\rho(t,d)\le L^{-2}\).
Both correction terms tend to zero as \(d\to\infty\).

This bound also handles sequences with oscillating dimensions.
Given any positive tolerance, choose a dimension threshold so that both corrections are sufficiently small above that threshold.
Then \(g_k\to0\) makes the displayed upper bound small for all sufficiently large \(k\) whose dimensions exceed the threshold.
Below it, the elementary upper comparison gives
\[
\rho(t_k,d_k)\le\frac{4d_k^2}{t_k},
\]
which tends uniformly to zero as \(t_k\to\infty\).
Hence
\[
\rho(t_k,d_k)\to0
\quad\Longleftrightarrow\quad
t_k\to\infty\ \text{and}\ g_k\to0.
\]
% lean: CausalSmith.Stat.LdpOptvalueUniformFrontier.rho_consistency_iff

\item \textit{Parametric comparison and the quadratic-resource regime.}
Suppose that fixed positive finite constants
\(c_{\mathrm{cmp}}\le C_{\mathrm{cmp}}\) satisfy, eventually,
\[
\frac{c_{\mathrm{cmp}}}{t_k}
\le\rho(t_k,d_k)
\le\frac{C_{\mathrm{cmp}}}{t_k}.
\]
Since \(\rho(t_k,d_k)\le1\) and
\(t_k\rho(t_k,d_k)\ge\min\{t_k,d_k\}/4\), these imply
\[
t_k\ge c_{\mathrm{cmp}},
\qquad
\min\{t_k,d_k\}\le4C_{\mathrm{cmp}}
\]
eventually.

Conversely, suppose that there are \(a_{\mathrm{par}}>0\) and
\(M_{\mathrm{par}}\in\mathbb R\) such that, eventually,
\[
t_k\ge a_{\mathrm{par}},
\qquad
\min\{t_k,d_k\}\le M_{\mathrm{par}}.
\]
Define
\[
B_{\mathrm{par}}=\max\{1,M_{\mathrm{par}}\}.
\]
The elementary lower comparison gives
\[
\frac{\min\{1,a_{\mathrm{par}}\}}{t_k}
\le\min\{1,t_k^{-1}\}\le\rho(t_k,d_k).
\]
For the upper comparison, split according to \(t_k\le d_k\) or \(d_k\le t_k\).
In the first case \(t_k\le B_{\mathrm{par}}\), so
\(\rho(t_k,d_k)\le1\le4B_{\mathrm{par}}^2/t_k\).
In the second case \(d_k\le B_{\mathrm{par}}\), and
\[
\rho(t_k,d_k)\le\frac{4d_k^2}{t_k}
\le\frac{4B_{\mathrm{par}}^2}{t_k}.
\]
This proves the asserted general parametric equivalence.
When \(t_k\to\infty\), the positive lower-resource condition holds eventually. Moreover,
\[
\min\{t_k,d_k\}\le M_{\mathrm{par}}
\]
eventually forces \(d_k\le M_{\mathrm{par}}\) eventually, since
\(t_k>M_{\mathrm{par}}\) eventually. Enlarging to an integer bound proves the bounded-dimension characterization; its converse follows from
\(\min\{t_k,d_k\}\le d_k\).

For the quadratic-resource regime, fix
\(0<a_{\mathrm{el}}\le b_{\mathrm{el}}\).
Choose an integer \(D_0\) large enough that, whenever \(d\ge D_0\),
\[
L\ge b_{\mathrm{el}}+1,\qquad
\log L\ge2\log b_{\mathrm{el}},\qquad
\log L\ge\log(e+a_{\mathrm{el}}^{-1}).
\]
For every \(t\) in the specified band, define
\[
u_{\mathrm{el}}=\frac{t}{d^2}\in[a_{\mathrm{el}},b_{\mathrm{el}}].
\]
Then \(t<d^2L\), and
\[
\log(e+L/u_{\mathrm{el}})
\ge\log(L/b_{\mathrm{el}})
\ge\frac12\log L,
\]
\[
\log(e+L/u_{\mathrm{el}})
\le\log\!\left(L(e+a_{\mathrm{el}}^{-1})\right)
\le2\log L.
\]
Since \(0\le(\log L)/L\le1\), applying the cap and squaring gives
\[
\frac14\left(\frac{\log L}{L}\right)^2
\le\rho(t,d)
\le4\left(\frac{\log L}{L}\right)^2.
\]
These bounds supply the required constants uniformly over the band.
% lean: CausalSmith.Stat.LdpOptvalueUniformFrontier.rate_resource_conclusions

\item \textit{Transfer to the actual minimax values.}
The upper bounds already proved make the minimax values finite, so their real values satisfy the same four inequalities. For a fixed class and an admissible sequence, define
\[
\begin{aligned}
\mathcal R_{\mathcal C,k}^{\mathrm{TV}}
&=\mathfrak R^{\mathrm{TV}}_{\mathcal C}
(\mathcal P_{d_k}^{\mathrm{pair}};n_k,d_k,\varepsilon_k),\\
\mathcal H_{\mathcal C,k}^{\mathrm{TV}}
&=\mathfrak H^{\mathrm{TV}}_{\mathcal C}
(\mathcal P_{d_k}^{\mathrm{pair}};n_k,d_k,\varepsilon_k).
\end{aligned}
\]
The common constants give
\[
c\rho(t_k,d_k)\le\mathcal R_{\mathcal C,k}^{\mathrm{TV}}
\le C_0\rho(t_k,d_k),
\]
\[
c\sqrt{\rho(t_k,d_k)}
\le\mathcal H_{\mathcal C,k}^{\mathrm{TV}}
\le C_0\sqrt{\rho(t_k,d_k)}.
\]
Squeezing nonnegative sequences therefore yields
\[
\mathcal R_{\mathcal C,k}^{\mathrm{TV}}\to0
\Longleftrightarrow\rho(t_k,d_k)\to0,
\]
\[
\mathcal H_{\mathcal C,k}^{\mathrm{TV}}\to0
\Longleftrightarrow\sqrt{\rho(t_k,d_k)}\to0
\Longleftrightarrow\rho(t_k,d_k)\to0.
\]
Combining these equivalences with the rate consistency criterion proves both minimax consistency assertions.

Two-sided comparison is symmetric and transitive: reversing
\(c\rho\le\mathcal R\le C_0\rho\) gives
\(\mathcal R/C_0\le\rho\le\mathcal R/c\).
Consequently,
\[
\mathcal R_{\mathcal C,k}^{\mathrm{TV}}\asymp t_k^{-1}
\Longleftrightarrow
\rho(t_k,d_k)\asymp t_k^{-1}.
\]
The general and diverging-resource characterizations thus transfer directly.

Finally, in the band from the preceding step, the common bounds imply
\[
\frac c4\left(\frac{\log L}{L}\right)^2
\le
\mathfrak R^{\mathrm{TV}}_{\mathcal C}
(\mathcal P_d^{\mathrm{pair}};n,d,\varepsilon)
\le
4C_0\left(\frac{\log L}{L}\right)^2,
\]
and
\[
\frac c2\left|\frac{\log L}{L}\right|
\le
\mathfrak H^{\mathrm{TV}}_{\mathcal C}
(\mathcal P_d^{\mathrm{pair}};n,d,\varepsilon)
\le
2C_0\left|\frac{\log L}{L}\right|.
\]
To give the risk and length bounds the coupled constants in the statement, define
\[
c_{\mathrm{el}}=\min\{c/4,c^2/4\},
\qquad
C_{\mathrm{el}}=\max\{4C_0,4C_0^2\}.
\]
Then \(0<c_{\mathrm{el}}\le C_{\mathrm{el}}<\infty\), and
\[
\sqrt{c_{\mathrm{el}}}\le c/2,
\qquad
2C_0\le\sqrt{C_{\mathrm{el}}}.
\]
The displayed risk and length inequalities therefore imply exactly the final two bounds, with \(c_{\mathrm{el}},C_{\mathrm{el}}\) as the band-dependent constants. This completes all the sequence conclusions.
% lean: CausalSmith.Stat.LdpOptvalueUniformFrontier.value_sequence_conclusions_of_bounds
\end{enumerate}
\end{proof}

\begin{proof}[Proof of \cref{obj:thm:full-simplex-tv-converse}]
The lower bounds on the full simplex follow from the paired-family lower bounds and model inclusion. We record both steps to make the interval scope explicit.

For every \(\theta\in[-1/2,1/2]^d\),
\[
P_\theta(j,s)=\frac{1+s\theta_j}{2d}\ge0,
\qquad
\sum_{s\in\{-1,1\}}P_\theta(j,s)=\frac1d.
\]
Summing over \(j\) shows that \(P_\theta\) is a probability law on \(\mathcal A_d\), hence
\[
\mathcal P_d^{\mathrm{pair}}\subseteq\Delta_{2d}.
\]
This is the inclusion in \cref{obj:def:tv-corollary-models}. % lean: CausalSmith.Stat.LdpOptvalueUniformFrontier.pairedFamily_subset_simplex

Fix \(\mathcal C\in\{\mathrm{NI},\mathrm{SI}\}\). By the converse part of \cref{obj:thm:paired-uniform-tv-frontier}, there is a universal constant \(c>0\), common to both protocol classes and decision criteria, such that
\[
\mathfrak R^{\mathrm{TV}}_{\mathcal C}
 (\mathcal P_d^{\mathrm{pair}};n,d,\varepsilon)
\ge c\,r_{\mathcal C}(n,d,\varepsilon)
\]
and
\[
\mathfrak H^{\mathrm{TV}}_{\mathcal C}
 (\mathcal P_d^{\mathrm{pair}};n,d,\varepsilon)
\ge c\,h_{\mathcal C}(n,d,\varepsilon).
\]
These are the paired lower bounds established under the same cited Bernstein premise stated in both theorem environments.

For squared-error risk, any admissible protocol-estimator pair on \(\Delta_{2d}\) is also admissible on the paired subfamily, and
\[
\sup_{\mathsf p\in\Delta_{2d}}
 \mathbb E_{\mathsf p,\mathsf K}
 \bigl[(\widehat D-D_{\mathrm{TV}}(\mathsf p))^2\bigr]
\ge
\sup_{\mathsf p\in\mathcal P_d^{\mathrm{pair}}}
 \mathbb E_{\mathsf p,\mathsf K}
 \bigl[(\widehat D-D_{\mathrm{TV}}(\mathsf p))^2\bigr].
\]
Taking the infimum over the common protocol and estimator class gives
\[
\mathfrak R^{\mathrm{TV}}_{\mathcal C}
 (\Delta_{2d};n,d,\varepsilon)
\ge
\mathfrak R^{\mathrm{TV}}_{\mathcal C}
 (\mathcal P_d^{\mathrm{pair}};n,d,\varepsilon)
\ge c\,r_{\mathcal C}(n,d,\varepsilon).
\]

For interval length, a connected closed interval in \([0,1]\) that has coverage at least \(0.90\) at every law in \(\Delta_{2d}\) has the same coverage on \(\mathcal P_d^{\mathrm{pair}}\). Its Borel endpoints, connectedness, and range therefore satisfy the paired decision requirements. Moreover,
\[
\sup_{\mathsf p\in\Delta_{2d}}
 \mathbb E_{\mathsf p,\mathsf K}\operatorname{length}J_{\mathrm{TV}}
\ge
\sup_{\mathsf p\in\mathcal P_d^{\mathrm{pair}}}
 \mathbb E_{\mathsf p,\mathsf K}\operatorname{length}J_{\mathrm{TV}}.
\]
Taking the infimum over globally simplex-honest protocol-interval pairs yields
\[
\mathfrak H^{\mathrm{TV}}_{\mathcal C}
 (\Delta_{2d};n,d,\varepsilon)
\ge
\mathfrak H^{\mathrm{TV}}_{\mathcal C}
 (\mathcal P_d^{\mathrm{pair}};n,d,\varepsilon)
\ge c\,h_{\mathcal C}(n,d,\varepsilon).
\]
Both inequalities use the same universal positive constant. % lean: CausalSmith.Stat.LdpOptvalueUniformFrontier.full_simplex_tv_converse
\end{proof}

\begin{proof}[Proof of \cref{obj:lem:bounded-mean-concentration}]
\begin{enumerate}
\item
Since \(m\ge1\) and a probability space is nonempty, the pointwise bound for \(H_1\) gives
\[
-b\le H_1\le b
\quad\Longrightarrow\quad
-b\le b
\quad\Longrightarrow\quad
0\le b.
\]
If \(b=0\), both asserted bounds follow from
\[
\Pr\left\{m^{-1}\sum_{i=1}^m(H_i-\mathbb EH_i)\ge x\right\}\le1,
\qquad
\Pr\left\{m^{-1}\sum_{i=1}^m(H_i-\mathbb EH_i)\le-x\right\}\le1.
\]
Henceforth suppose \(b>0\), and define the variance proxy
\[
c_b:=\left(\frac{|b-(-b)|}{2}\right)^2=b^2>0.
\]
% lean: CausalSmith.Stat.LdpOptvalueUniformFrontier.boundedMeanConcentration_proved

\item
We first establish an upper-tail bound for any mutually independent, measurable real random variables \(H_1^\star,\ldots,H_m^\star\) on the given probability space satisfying
\[
-b\le H_i^\star\le b,
\qquad i=1,\ldots,m.
\]
Define their centered versions and interval endpoints by
\[
H_i^\circ:=H_i^\star-\mathbb EH_i^\star,
\qquad
\ell_i^\circ:=-b-\mathbb EH_i^\star,
\qquad
u_i^\circ:=b-\mathbb EH_i^\star.
\]
Boundedness ensures integrability, and therefore
\[
\mathbb EH_i^\circ=0,
\qquad
\ell_i^\circ\le H_i^\circ\le u_i^\circ,
\qquad
u_i^\circ-\ell_i^\circ=2b.
\]
The centered variables remain mutually independent because each is a measurable function of its corresponding original variable.

For each \(i\), define the moment-generating function and its logarithm on all of \(\mathbb R\) by
\[
\Phi_i(\upsilon):=\mathbb E\exp(\upsilon H_i^\circ),
\qquad
\psi_i(\upsilon):=\log\Phi_i(\upsilon),
\qquad \upsilon\in\mathbb R.
\]
Boundedness makes \(\Phi_i\) finite and strictly positive and permits differentiation under the expectation. Define expectation under the exponential tilt by
\[
\mathbb E_{\upsilon}^{(i)}[\cdot]
:=
\frac{\mathbb E[\cdot\,\exp(\upsilon H_i^\circ)]}{\Phi_i(\upsilon)},
\qquad
\bar h_{i,\upsilon}:=\mathbb E_{\upsilon}^{(i)}H_i^\circ.
\]
This is expectation under a probability measure absolutely continuous with respect to the original measure. Direct differentiation gives
\[
\psi_i(0)=0,
\qquad
\psi_i'(0)=\mathbb EH_i^\circ=0,
\qquad
\psi_i''(\upsilon)
=
\mathbb E_{\upsilon}^{(i)}[(H_i^\circ)^2]
-\bar h_{i,\upsilon}^{\,2}.
\]

The interval bound persists under the tilt. Consequently,
\[
\begin{aligned}
0
&\le
\mathbb E_{\upsilon}^{(i)}
\bigl[(u_i^\circ-H_i^\circ)(H_i^\circ-\ell_i^\circ)\bigr]\\
&=
-\mathbb E_{\upsilon}^{(i)}[(H_i^\circ)^2]
+(u_i^\circ+\ell_i^\circ)\bar h_{i,\upsilon}
-u_i^\circ\ell_i^\circ.
\end{aligned}
\]
Subtracting the squared tilted mean yields
\[
\begin{aligned}
\psi_i''(\upsilon)
&\le
(u_i^\circ-\bar h_{i,\upsilon})
(\bar h_{i,\upsilon}-\ell_i^\circ)\\
&=
\left(\frac{u_i^\circ-\ell_i^\circ}{2}\right)^2
-
\left(\bar h_{i,\upsilon}
-\frac{u_i^\circ+\ell_i^\circ}{2}\right)^2
\le c_b.
\end{aligned}
\]
For \(\upsilon>0\), Taylor's theorem supplies an intermediate point
\[
\upsilon_{i,\mathrm{mid}}\in(0,\upsilon),
\qquad
\psi_i(\upsilon)
=
\frac{\upsilon^2}{2}\psi_i''(\upsilon_{i,\mathrm{mid}})
\le\frac{c_b\upsilon^2}{2}.
\]
Thus \(\Phi_i(\upsilon)\le\exp(c_b\upsilon^2/2)\) for positive \(\upsilon\); at zero both sides equal one. For negative \(\upsilon\), apply the same positive-parameter calculation to the centered variable \(-H_i^\circ\), whose interval has the same width:
\[
-u_i^\circ\le-H_i^\circ\le-\ell_i^\circ,
\qquad
\mathbb E[-H_i^\circ]=0.
\]
It follows that
\[
\Phi_i(\upsilon)
=
\mathbb E\exp\bigl((-\upsilon)(-H_i^\circ)\bigr)
\le
\exp\left(\frac{c_b(-\upsilon)^2}{2}\right).
\]
We have therefore established
\[
\mathbb E\exp(\upsilon H_i^\circ)
\le\exp\left(\frac{c_b\upsilon^2}{2}\right),
\qquad \upsilon\in\mathbb R,\quad i=1,\ldots,m.
\]
% lean: ProbabilityTheory.hasSubgaussianMGF_of_mem_Icc

\item
Define the centered sum
\[
S^\star:=\sum_{i=1}^m H_i^\circ.
\]
Independence gives the factorization
\[
\begin{aligned}
\mathbb E\exp(\upsilon S^\star)
&=\prod_{i=1}^m\mathbb E\exp(\upsilon H_i^\circ)\\
&\le
\prod_{i=1}^m\exp\left(\frac{c_b\upsilon^2}{2}\right)
=
\exp\left(\frac{mc_b\upsilon^2}{2}\right),
\qquad \upsilon\in\mathbb R.
\end{aligned}
\]
Indeed, adjoining one centered variable to a partial sum multiplies their moment-generating functions, since that variable is independent of the preceding sum; induction gives the displayed formula.

Since \(mx\ge0\), for every \(\upsilon\ge0\) exponential Markov's inequality gives
\[
\begin{aligned}
\Pr\{S^\star\ge mx\}
&\le
\exp(-\upsilon mx)\,\mathbb E\exp(\upsilon S^\star)\\
&\le
\exp\left(-\upsilon mx+\frac{mc_b\upsilon^2}{2}\right).
\end{aligned}
\]
Choose
\[
\upsilon_{\mathrm{opt}}:=\frac{mx}{mc_b}=\frac{x}{c_b}\ge0.
\]
This choice is defined also when \(x=0\), and substitution yields
\[
\Pr\{S^\star\ge mx\}
\le
\exp\left(-\frac{(mx)^2}{2mc_b}\right).
\]
% lean: ProbabilityTheory.HasSubgaussianMGF.measure_sum_ge_le_of_iIndepFun

\item
Because \(m>0\), multiplying by \(m\) gives the event identity
\[
\left\{
x\le m^{-1}\sum_{i=1}^m
(H_i^\star-\mathbb EH_i^\star)
\right\}
=
\{mx\le S^\star\}.
\]
Using \(c_b=b^2\), the exponent simplifies to
\[
-\frac{(mx)^2}{2mc_b}
=
-\frac{mx^2}{2b^2}.
\]
Hence the upper-tail calculation establishes
\[
\Pr\left\{
m^{-1}\sum_{i=1}^m(H_i^\star-\mathbb EH_i^\star)\ge x
\right\}
\le
\exp\left(-\frac{mx^2}{2b^2}\right).
\]
Taking \(H_i^\star=H_i\) proves the asserted upper-tail bound.
% lean: CausalSmith.Stat.LdpOptvalueUniformFrontier.boundedMeanConcentration_proved.upper

\item
For the lower tail, define
\[
H_i^-:=-H_i,
\qquad i=1,\ldots,m.
\]
These are mutually independent and measurable, and negating the original bounds gives
\[
-b\le H_i^-\le b.
\]
Linearity of expectation and finite summation gives
\[
\mathbb EH_i^-=-\mathbb EH_i,
\qquad
\sum_{i=1}^m(H_i^--\mathbb EH_i^-)
=
-\sum_{i=1}^m(H_i-\mathbb EH_i).
\]
Therefore,
\[
\left\{
m^{-1}\sum_{i=1}^m(H_i-\mathbb EH_i)\le-x
\right\}
=
\left\{
m^{-1}\sum_{i=1}^m(H_i^--\mathbb EH_i^-)\ge x
\right\}.
\]
Applying the upper-tail bound just established to \(H_1^-,\ldots,H_m^-\) yields
\[
\Pr\left\{
m^{-1}\sum_{i=1}^m(H_i-\mathbb EH_i)\le-x
\right\}
\le
\exp\left(-\frac{mx^2}{2b^2}\right),
\]
which completes the proof.
% lean: CausalSmith.Stat.LdpOptvalueUniformFrontier.boundedMeanConcentration_proved
\end{enumerate}
\end{proof}

\begin{proof}[Proof of \cref{obj:lem:cai-low-chebyshev-approximation}]
\begin{enumerate}
\item Fix an even integer \(D\ge2\), and set
\[
\ell_{\mathrm{tr}}:=D/2\in\mathbb N,\qquad
\ell_{\mathrm{tr}}\ge1,\qquad D=2\ell_{\mathrm{tr}}.
\]
For the truncation argument, extend the polynomial family to degree zero by
\[
q_0(x):=\frac2\pi.
\]
Thus, for every \(\ell\in\mathbb N_0\),
\[
q_{2\ell}(x)
=\frac2\pi+\frac4\pi\sum_{\ell'=1}^{\ell}
 \frac{(-1)^{\ell'+1}\mathcal T_{2\ell'}(x)}
 {4(\ell')^2-1},
\qquad q_D=q_{2\ell_{\mathrm{tr}}}.
\]
This is precisely the one-based truncation of the even Chebyshev modes; the sum is empty when \(\ell=0\).
% lean: CausalSmith.Stat.LdpOptvalueUniformFrontier.chebyshevAbsPoly_eq_library

\item Define the continuous, \(\pi\)-periodic function \(\varphi_{\sin}:\mathbb R\to\mathbb R\) by
\[
\varphi_{\sin}(\vartheta)=\sin\vartheta
\quad(0\le\vartheta\le\pi),\qquad
\varphi_{\sin}(\vartheta+\pi)=\varphi_{\sin}(\vartheta)
\quad(\vartheta\in\mathbb R).
\]
The endpoint values agree, so the periodic extension is continuous. With \(\mathrm i:=\sqrt{-1}\), define its Fourier coefficients by
\[
\widehat\varphi_{\sin}(\ell)
:=\frac1\pi\int_0^\pi
 \sin\vartheta\,e^{-2\mathrm i\ell\vartheta}\,d\vartheta,
\qquad \ell\in\mathbb Z.
\]
For each such \(\ell\), \(1-4\ell^2\ne0\), and the function
\[
\Phi_\ell(\vartheta)
:=e^{-2\mathrm i\ell\vartheta}
 \frac{-2\mathrm i\ell\sin\vartheta-\cos\vartheta}
 {1-4\ell^2},
\qquad \vartheta\in\mathbb R,
\]
satisfies
\[
\Phi_\ell'(\vartheta)
=e^{-2\mathrm i\ell\vartheta}\sin\vartheta.
\]
Since \(e^{-2\mathrm i\ell\pi}=1\), evaluation at the endpoints gives
\[
\widehat\varphi_{\sin}(\ell)
=\frac{\Phi_\ell(\pi)-\Phi_\ell(0)}{\pi}
=\frac{2}{\pi(1-4\ell^2)}
=-\frac{2}{\pi(4\ell^2-1)}.
\]
In particular, \(\widehat\varphi_{\sin}(0)=2/\pi\).
% lean: Causalean.Mathlib.Analysis.AbsoluteValueMomentPriorDuality.fourierCoeff_sinCircle

\item For every integer \(\ell\ge1\),
\[
\bigl|\widehat\varphi_{\sin}(\ell)\bigr|
=\frac{2}{\pi(4\ell^2-1)}
\le\frac1{\ell^2}.
\]
Indeed, \(\ell^2\ge1\) and \(\pi\ge2\) imply
\[
2\ell^2\le\pi(4\ell^2-1).
\]
The coefficients at \(\ell\) and \(-\ell\) agree. Comparison with the convergent series \(\sum_{\ell=1}^\infty\ell^{-2}\) therefore yields
\[
\sum_{\ell\in\mathbb Z}
 \bigl|\widehat\varphi_{\sin}(\ell)\bigr|<\infty.
\]
The Fourier reconstruction theorem for a continuous periodic function with absolutely summable coefficients now gives
\[
\varphi_{\sin}(\vartheta)
=\sum_{\ell\in\mathbb Z}
 \widehat\varphi_{\sin}(\ell)e^{2\mathrm i\ell\vartheta},
\qquad \vartheta\in\mathbb R.
\]
% lean: Causalean.Mathlib.Analysis.AbsoluteValueMomentPriorDuality.summable_fourierCoeff_sinCircle

\item For \(0\le\vartheta\le\pi\), periodicity and the sine addition identities give
\[
\varphi_{\sin}(\vartheta+\pi/2)
=
\begin{cases}
\cos\vartheta,&0\le\vartheta\le\pi/2,\\
-\cos\vartheta,&\pi/2<\vartheta\le\pi,
\end{cases}
=|\cos\vartheta|.
\]
The defining Chebyshev recurrence and the cosine addition formula also give
\[
\mathcal T_\ell(\cos\vartheta)=\cos(\ell\vartheta),
\qquad \ell\in\mathbb N_0,\quad \vartheta\in\mathbb R.
\]
Consequently, for \(\ell\ge1\), the paired Fourier terms at the shifted argument satisfy
\[
\begin{aligned}
&\widehat\varphi_{\sin}(\ell)
 e^{2\mathrm i\ell(\vartheta+\pi/2)}
+\widehat\varphi_{\sin}(-\ell)
 e^{-2\mathrm i\ell(\vartheta+\pi/2)}
\\
&\qquad
=-\frac{4\cos(2\ell\vartheta+\ell\pi)}
 {\pi(4\ell^2-1)}
=\frac4\pi
 \frac{(-1)^{\ell+1}\mathcal T_{2\ell}(\cos\vartheta)}
 {4\ell^2-1}.
\end{aligned}
\]
For \(x\in[-1,1]\), define
\[
\vartheta_x:=\arccos x\in[0,\pi],
\qquad \cos\vartheta_x=x.
\]
Pairing the positive and negative frequencies in the absolutely convergent Fourier series and taking real parts therefore yields
\[
|x|=\frac2\pi+\frac4\pi\sum_{\ell=1}^{\infty}
 \frac{(-1)^{\ell+1}\mathcal T_{2\ell}(x)}
 {4\ell^2-1}.
\]
% lean: Causalean.Mathlib.Analysis.AbsoluteValueMomentPriorDuality.hasSum_absChebTerm_real

\item Splitting this series at \(\ell_{\mathrm{tr}}\) gives
\[
|x|-q_D(x)
=\frac4\pi\sum_{\ell=\ell_{\mathrm{tr}}+1}^{\infty}
 \frac{(-1)^{\ell+1}\mathcal T_{2\ell}(x)}
 {4\ell^2-1}.
\]
The scalar majorant has an exact telescoping sum, because
\[
\frac1{4\ell^2-1}
=\frac12\left(\frac1{2\ell-1}-\frac1{2\ell+1}\right)
\qquad(\ell\ge1).
\]
In particular, for every integer \(\ell_{\mathrm{end}}\ge\ell_{\mathrm{tr}}\),
\[
\sum_{\ell=\ell_{\mathrm{tr}}+1}^{\ell_{\mathrm{end}}}
 \frac1{4\ell^2-1}
=\frac12\left(
 \frac1{2\ell_{\mathrm{tr}}+1}
 -\frac1{2\ell_{\mathrm{end}}+1}\right),
\]
including the empty sum when \(\ell_{\mathrm{end}}=\ell_{\mathrm{tr}}\). Passing to the limit gives
\[
\sum_{\ell=\ell_{\mathrm{tr}}+1}^{\infty}
 \frac1{4\ell^2-1}
=\frac1{2(2\ell_{\mathrm{tr}}+1)}.
\]
Moreover, the cosine identity implies
\[
|\mathcal T_{2\ell}(x)|
=|\cos(2\ell\vartheta_x)|\le1.
\]
Thus the residual series is absolutely convergent, and
\[
\begin{aligned}
\bigl|q_D(x)-|x|\bigr|
&\le\frac4\pi
 \sum_{\ell=\ell_{\mathrm{tr}}+1}^{\infty}
 \frac{|\mathcal T_{2\ell}(x)|}{4\ell^2-1}\\
&\le\frac4\pi
 \sum_{\ell=\ell_{\mathrm{tr}}+1}^{\infty}
 \frac1{4\ell^2-1}
=\frac2{\pi(2\ell_{\mathrm{tr}}+1)}
=\frac2{\pi(D+1)}.
\end{aligned}
\]
% lean: Causalean.Mathlib.Analysis.AbsoluteValueMomentPriorDuality.absChebPoly_error_le

\item For real polynomials
\[
\mathcal B(x)=\sum_{\ell\ge0}\beta_\ell x^\ell,
\qquad
\mathcal E(x)=\sum_{\ell\ge0}\gamma_\ell x^\ell,
\]
where the real coefficient sequences have finite support, define
\[
\|\mathcal B\|_{\mathrm{coef},1}
:=\sum_{\ell\ge0}|\beta_\ell|,
\qquad
\|\mathcal E\|_{\mathrm{coef},1}
:=\sum_{\ell\ge0}|\gamma_\ell|.
\]
Coefficientwise triangle inequalities give
\[
\|\mathcal B\pm\mathcal E\|_{\mathrm{coef},1}
\le
\|\mathcal B\|_{\mathrm{coef},1}
+\|\mathcal E\|_{\mathrm{coef},1}.
\]
For every real scalar \(\zeta\), the same definition gives
\[
\|0\|_{\mathrm{coef},1}=0,\qquad
\|\zeta\|_{\mathrm{coef},1}=|\zeta|,\qquad
\|\zeta\mathcal B\|_{\mathrm{coef},1}
=|\zeta|\|\mathcal B\|_{\mathrm{coef},1},\qquad
\|x\mathcal B\|_{\mathrm{coef},1}
=\|\mathcal B\|_{\mathrm{coef},1}.
\]
Here the last identity follows by shifting the coefficient indices. Finally, expanding the product into monomials and applying the finite-sum triangle inequality yields
\[
\begin{aligned}
\mathcal B(x)\mathcal E(x)
&=\sum_{\ell,\ell'\ge0}
 \beta_\ell\gamma_{\ell'}x^{\ell+\ell'},\\
\|\mathcal B\mathcal E\|_{\mathrm{coef},1}
&\le\sum_{\ell,\ell'\ge0}
 |\beta_\ell\gamma_{\ell'}|
=\|\mathcal B\|_{\mathrm{coef},1}
 \|\mathcal E\|_{\mathrm{coef},1}.
\end{aligned}
\]
All sums here have finitely many nonzero terms.
% lean: Causalean.Mathlib.Analysis.Approximation.Chebyshev.polynomialCoeffOneNorm_mul_le

\item Define
\[
\chi:=1+\sqrt2>0,\qquad \chi^2=2\chi+1.
\]
We claim that
\[
\|\mathcal T_\ell\|_{\mathrm{coef},1}\le\chi^\ell
\qquad(\ell\in\mathbb N_0).
\]
The two initial cases are
\[
\|\mathcal T_0\|_{\mathrm{coef},1}=1=\chi^0,
\qquad
\|\mathcal T_1\|_{\mathrm{coef},1}=1\le\chi.
\]
For the induction step, the Chebyshev recurrence and the norm identities just established give
\[
\begin{aligned}
\|\mathcal T_{\ell+2}\|_{\mathrm{coef},1}
&=\|2x\mathcal T_{\ell+1}-\mathcal T_\ell\|_{\mathrm{coef},1}\\
&\le2\|\mathcal T_{\ell+1}\|_{\mathrm{coef},1}
 +\|\mathcal T_\ell\|_{\mathrm{coef},1}\\
&\le2\chi^{\ell+1}+\chi^\ell
=\chi^\ell(2\chi+1)
=\chi^{\ell+2}.
\end{aligned}
\]
This proves the claim by two-step induction.
% lean: Causalean.Mathlib.Analysis.Approximation.Chebyshev.polynomialCoeffOneNorm_chebyshev_le

\item For every \(\ell\in\mathbb N_0\),
\[
4(\ell+1)^2-1\ge3,\qquad
\frac4\pi\le\frac43,
\]
where the second inequality uses \(\pi>3\). Hence
\[
\left|
 \frac{(4/\pi)(-1)^\ell}{4(\ell+1)^2-1}
\right|
=\frac{4/\pi}{4(\ell+1)^2-1}
\le\frac12.
\]
Also, \(2\le(3/2)^2\) implies \(\sqrt2\le3/2\), and therefore
\[
\chi^2=3+2\sqrt2\le6.
\]
The product bound and the Chebyshev norm estimate now imply
\[
\begin{aligned}
\left\|
 \frac{(4/\pi)(-1)^\ell}{4(\ell+1)^2-1}
 \mathcal T_{2(\ell+1)}
\right\|_{\mathrm{coef},1}
&\le\frac12\chi^{2(\ell+1)}\\
&\le\frac12\,6^{\ell+1}
=3\cdot6^\ell
\le3\cdot8^\ell.
\end{aligned}
\]
% lean: Causalean.Mathlib.Analysis.Approximation.Chebyshev.AbsoluteValue.absChebPoly_coeffOneNorm_le

\item We next prove, by induction on \(\ell\in\mathbb N_0\), that
\[
\|q_{2\ell}\|_{\mathrm{coef},1}\le8^\ell.
\]
At \(\ell=0\),
\[
\|q_0\|_{\mathrm{coef},1}=\frac2\pi\le1=8^0.
\]
The truncations satisfy
\[
q_{2(\ell+1)}
=q_{2\ell}
+\frac{(4/\pi)(-1)^\ell}{4(\ell+1)^2-1}
 \mathcal T_{2(\ell+1)}.
\]
Thus the induction hypothesis and the preceding mode bound give
\[
\begin{aligned}
\|q_{2(\ell+1)}\|_{\mathrm{coef},1}
&\le\|q_{2\ell}\|_{\mathrm{coef},1}
+\left\|
 \frac{(4/\pi)(-1)^\ell}{4(\ell+1)^2-1}
 \mathcal T_{2(\ell+1)}
 \right\|_{\mathrm{coef},1}\\
&\le8^\ell+3\cdot8^\ell
\le8^{\ell+1}.
\end{aligned}
\]
In particular,
\[
\|q_D\|_{\mathrm{coef},1}
\le8^{\ell_{\mathrm{tr}}}
=2^{3\ell_{\mathrm{tr}}}
=2^{3D/2}.
\]
% lean: Causalean.Mathlib.Analysis.Approximation.Chebyshev.AbsoluteValue.absChebPoly_coeffOneNorm_le

\item Fix an integer \(v\) with \(0\le v\le D\). If \(a_{Dv}=0\), then
\[
|a_{Dv}|=0\le2^{3D/2}.
\]
Otherwise, its absolute value is one of the nonnegative summands in the coefficient norm. Using the stated monomial expansion of \(q_D\), we obtain
\[
|a_{Dv}|
\le\sum_{\ell=0}^{D}|a_{D\ell}|
=\|q_D\|_{\mathrm{coef},1}
\le2^{3D/2}.
\]
Together with the approximation bound, this proves both conclusions.
% lean: Causalean.Mathlib.Analysis.Approximation.Chebyshev.AbsoluteValue.absChebPoly_halfDegree_coeff_le
\end{enumerate}
\end{proof}

\begin{proof}[Proof of \cref{obj:lem:cai-low-moment-duality}]
Fix a positive even integer \(K\).

\begin{enumerate}
\item Define the approximation space and target function by
\[
I_{\mathrm{appr}}=[-1,1],\qquad
\mathcal C_{\mathrm{appr}}=C(I_{\mathrm{appr}},\mathbb R),\qquad
f_{\mathrm{abs}}(x)=|x|,
\]
and equip \(\mathcal C_{\mathrm{appr}}\) with the uniform norm. Let
\[
\mathcal E_K
=\{p|_{I_{\mathrm{appr}}}:p\in\mathbb R[x],\ \deg p\le K\}.
\]
For every real polynomial \(p\),
\[
\|f_{\mathrm{abs}}-p|_{I_{\mathrm{appr}}}\|_\infty
=\sup_{|x|\le1}\bigl||x|-p(x)\bigr|.
\]
Thus \(e_K\) is the distance from \(f_{\mathrm{abs}}\) to \(\mathcal E_K\).

We first verify attainment of this distance. Define
\[
R_{\mathrm{ball}}=2\|f_{\mathrm{abs}}\|_\infty+1,\qquad
\mathcal B_K
=\{h\in\mathcal E_K:\|h\|_\infty\le R_{\mathrm{ball}}\}.
\]
The space \(\mathcal E_K\) is finite dimensional, so \(\mathcal B_K\) is compact. The continuous function
\(h\mapsto\|f_{\mathrm{abs}}-h\|_\infty\) therefore has a minimizer
\(h_{\mathrm{min}}\in\mathcal B_K\). Since \(0\in\mathcal B_K\),
\[
\|f_{\mathrm{abs}}-h_{\mathrm{min}}\|_\infty
\le\|f_{\mathrm{abs}}\|_\infty.
\]
For \(h\in\mathcal E_K\setminus\mathcal B_K\), the reverse triangle inequality gives
\[
\|f_{\mathrm{abs}}-h\|_\infty
\ge\|h\|_\infty-\|f_{\mathrm{abs}}\|_\infty
>\|f_{\mathrm{abs}}\|_\infty+1.
\]
Consequently \(h_{\mathrm{min}}\) minimizes over all of \(\mathcal E_K\), and
\[
\|f_{\mathrm{abs}}-h_{\mathrm{min}}\|_\infty=e_K.
\]
% lean: Causalean.Mathlib.Analysis.AbsoluteValueMomentPriorDuality.exists_bestPolynomialAbs

\item Let the quotient map and the target's quotient class be
\[
\pi_{\mathrm{quot}}:\mathcal C_{\mathrm{appr}}
\longrightarrow\mathcal C_{\mathrm{appr}}/\mathcal E_K,
\qquad
\bar f_{\mathrm{abs}}=\pi_{\mathrm{quot}}(f_{\mathrm{abs}}).
\]
The quotient norm is the distance to the subspace. The minimizer above gives its upper bound, while the definition of \(e_K\) gives its lower bound:
\[
\|\bar f_{\mathrm{abs}}\|
=\inf_{h\in\mathcal E_K}\|f_{\mathrm{abs}}-h\|_\infty
=e_K.
\]
By the norming-functional form of Hahn--Banach, there is a continuous real linear functional
\[
\gamma_K:\mathcal C_{\mathrm{appr}}/\mathcal E_K\longrightarrow\mathbb R,
\qquad
\|\gamma_K\|\le1,\qquad
\gamma_K(\bar f_{\mathrm{abs}})=\|\bar f_{\mathrm{abs}}\|.
\]
This also covers a zero quotient class, for which the zero functional suffices. Define
\[
\Lambda_K=\gamma_K\circ\pi_{\mathrm{quot}}.
\]
Since the quotient map is contractive and vanishes on \(\mathcal E_K\),
\[
|\Lambda_K(h)|\le\|h\|_\infty
\quad(h\in\mathcal C_{\mathrm{appr}}),\qquad
\Lambda_K(1)=0,\qquad
\Lambda_K(f_{\mathrm{abs}})=e_K.
\]
In particular, with
\[
h_\ell(x)=x^\ell
\qquad(x\in I_{\mathrm{appr}},\ 0\le\ell\le K),
\]
we have
\[
\Lambda_K(h_\ell)=0.
\]
% lean: Causalean.Mathlib.Analysis.AbsoluteValueMomentPriorDuality.exists_dual_abs_certificate

\item We construct two positive linear functionals of mass one half whose difference is \(\Lambda_K\). On
\(\mathcal C_{\mathrm{appr}}\times\mathbb R\), define the subspace, its linear functional, and a cone by
\[
\begin{aligned}
\mathcal D_K
&=\{(c_{\mathrm{const}}1,s_{\mathrm{aux}}):
       c_{\mathrm{const}},s_{\mathrm{aux}}\in\mathbb R\},\\
\beta_K(c_{\mathrm{const}}1,s_{\mathrm{aux}})
&=\frac{c_{\mathrm{const}}}{2}+s_{\mathrm{aux}},\\
\mathcal S_K
&=\{(h_++h_-,-\Lambda_K(h_-)):
       h_+,h_-\in\mathcal C_{\mathrm{appr}},\
       h_+\ge0,\ h_-\ge0\}.
\end{aligned}
\]
The representation of an element of \(\mathcal D_K\) is unique because
\(I_{\mathrm{appr}}\) is nonempty. The set \(\mathcal S_K\) is closed under addition and multiplication by nonnegative scalars.

Suppose
\[
(c_{\mathrm{const}}1,s_{\mathrm{aux}})
=(h_++h_-,-\Lambda_K(h_-))\in\mathcal D_K\cap\mathcal S_K.
\]
Then
\[
c_{\mathrm{const}}\ge0,\qquad
0\le h_-(x)\le c_{\mathrm{const}},
\]
so the function
\[
h_{\mathrm{cent}}=2h_--c_{\mathrm{const}}1
\]
satisfies
\[
\|h_{\mathrm{cent}}\|_\infty\le c_{\mathrm{const}},\qquad
2\Lambda_K(h_-)
=\Lambda_K(h_{\mathrm{cent}})
\le|\Lambda_K(h_{\mathrm{cent}})|
\le c_{\mathrm{const}}.
\]
It follows that
\[
\beta_K(c_{\mathrm{const}}1,s_{\mathrm{aux}})
=\frac{c_{\mathrm{const}}}{2}-\Lambda_K(h_-)\ge0.
\]
Moreover, for every
\((h_{\mathrm{test}},s_{\mathrm{test}})
\in\mathcal C_{\mathrm{appr}}\times\mathbb R\),
\[
(\|h_{\mathrm{test}}\|_\infty1,-s_{\mathrm{test}})
\in\mathcal D_K,
\]
and
\[
(\|h_{\mathrm{test}}\|_\infty1,-s_{\mathrm{test}})
+(h_{\mathrm{test}},s_{\mathrm{test}})
=(\|h_{\mathrm{test}}\|_\infty1+h_{\mathrm{test}},0)
\in\mathcal S_K,
\]
because the first coordinate is nonnegative. The ordered linear extension theorem therefore supplies a linear functional
\[
\Gamma_K:\mathcal C_{\mathrm{appr}}\times\mathbb R\longrightarrow\mathbb R,
\qquad
\Gamma_K|_{\mathcal D_K}=\beta_K,\qquad
\Gamma_K\ge0\ \text{on }\mathcal S_K.
\]
Define, for \(h\in\mathcal C_{\mathrm{appr}}\),
\[
\Phi_+(h)=\Gamma_K(h,0),\qquad
\Phi_-(h)=\Phi_+(h)-\Lambda_K(h).
\]
For \(h\ge0\), both \((h,0)\) and \((h,-\Lambda_K(h))\) belong to \(\mathcal S_K\). Since \(\Gamma_K(0,s)=s\) for every real \(s\),
\[
\Phi_+(h)\ge0,\qquad
\Phi_-(h)=\Gamma_K(h,-\Lambda_K(h))\ge0.
\]
Finally,
\[
\Phi_+(1)=\frac12,\qquad
\Phi_-(1)=\frac12,\qquad
\Phi_+(h)-\Phi_-(h)=\Lambda_K(h).
\]
% lean: Causalean.Mathlib.Analysis.AbsoluteValueMomentPriorDuality.exists_positive_half_split

\item Riesz representation on the compact interval gives finite positive Borel measures \(\mu_+,\mu_-\) on \(I_{\mathrm{appr}}\), characterized by
\[
\int_{I_{\mathrm{appr}}}h\,d\mu_+=\Phi_+(h),\qquad
\int_{I_{\mathrm{appr}}}h\,d\mu_-=\Phi_-(h)
\quad(h\in\mathcal C_{\mathrm{appr}}).
\]
Applying these identities to the constant function yields
\[
\mu_+(I_{\mathrm{appr}})=\mu_-(I_{\mathrm{appr}})=\frac12.
\]
Define the inclusion and reflection maps, and the symmetrized measures on \(\mathbb R\), by
\[
\iota:I_{\mathrm{appr}}\longrightarrow\mathbb R,\quad \iota(x)=x,
\qquad
r_{\mathrm{refl}}:\mathbb R\longrightarrow\mathbb R,\quad
r_{\mathrm{refl}}(x)=-x,
\]
\[
\zeta_\pm
=\frac12\left(\iota_\#\mu_\pm
 +(r_{\mathrm{refl}}\circ\iota)_\#\mu_\pm\right).
\]
Here \(\#\) denotes pushforward. Both pushforwards have mass one half and are supported on \(I_{\mathrm{appr}}\); reflection interchanges them. Hence
\[
\zeta_\pm(\mathbb R)=\frac12,\qquad
\zeta_\pm(\mathbb R\setminus I_{\mathrm{appr}})=0,\qquad
(r_{\mathrm{refl}})_\#\zeta_\pm=\zeta_\pm.
\]
For every continuous \(h_{\mathrm{test}}:\mathbb R\to\mathbb R\), the pushforward integral formula gives
\[
\int_{\mathbb R}h_{\mathrm{test}}\,d\zeta_\pm
=\frac12\left(
\int_{I_{\mathrm{appr}}}h_{\mathrm{test}}(x)\,d\mu_\pm(x)
+\int_{I_{\mathrm{appr}}}h_{\mathrm{test}}(-x)\,d\mu_\pm(x)
\right).
\]
All these integrals exist because the measures are finite and the functions are bounded on the relevant compact interval.
% lean: Causalean.Mathlib.Analysis.AbsoluteValueMomentPriorDuality.symmPush_mass
% lean: Causalean.Mathlib.Analysis.AbsoluteValueMomentPriorDuality.symmPush_supported
% lean: Causalean.Mathlib.Analysis.AbsoluteValueMomentPriorDuality.symmPush_symmetric
% lean: Causalean.Mathlib.Analysis.AbsoluteValueMomentPriorDuality.integral_symmPush

\item For \(0\le\ell\le K\), define the reflected monomial by
\[
h_{\ell,\mathrm{refl}}(x)=(-x)^\ell=(-1)^\ell h_\ell(x)
\qquad(x\in I_{\mathrm{appr}}).
\]
Linearity and polynomial annihilation give
\[
\Lambda_K(h_\ell)=0,\qquad
\Lambda_K(h_{\ell,\mathrm{refl}})
=(-1)^\ell\Lambda_K(h_\ell)=0.
\]
Combining the integral formula with the difference of the representing functionals therefore yields
\[
\begin{aligned}
\int_{\mathbb R}x^\ell\,d\zeta_+
-\int_{\mathbb R}x^\ell\,d\zeta_-
&=\frac12\left(
\Lambda_K(h_\ell)+\Lambda_K(h_{\ell,\mathrm{refl}})
\right)\\
&=0.
\end{aligned}
\]
For the absolute-value function, reflection preserves the function itself:
\[
|-x|=|x|.
\]
Consequently,
\[
\int_{\mathbb R}|x|\,d\zeta_+
-\int_{\mathbb R}|x|\,d\zeta_-
=\frac12\left(
\Lambda_K(f_{\mathrm{abs}})+\Lambda_K(f_{\mathrm{abs}})
\right)
=e_K.
\]
% lean: Causalean.Mathlib.Analysis.AbsoluteValueMomentPriorDuality.exists_absExtremalDecomposition

\item Normalize the two measures by defining
\[
\xi_0=2\zeta_-,\qquad \xi_1=2\zeta_+.
\]
Their masses are one, and multiplication by two preserves their support and reflection invariance:
\[
\xi_0(\mathbb R)=\xi_1(\mathbb R)=1,\qquad
\xi_0(\mathbb R\setminus[-1,1])
=\xi_1(\mathbb R\setminus[-1,1])=0,
\]
\[
(r_{\mathrm{refl}})_\#\xi_0=\xi_0,\qquad
(r_{\mathrm{refl}})_\#\xi_1=\xi_1.
\]
Scaling the moment identities gives, for every integer \(0\le\ell\le K\),
\[
\int x^\ell\,d\xi_1(x)-\int x^\ell\,d\xi_0(x)
=2\left(\int x^\ell\,d\zeta_+(x)-\int x^\ell\,d\zeta_-(x)\right)
=0.
\]
Likewise, scaling the absolute-moment identity gives
\[
\int |x|\,d\xi_1(x)-\int |x|\,d\xi_0(x)=2e_K.
\]
The approximation error used throughout agrees with the stated one because
\[
\inf_{\deg p\le K}\sup_{|x|\le1}\bigl||x|-p(x)\bigr|
=
\inf_{\deg p\le K}\sup_{|x|\le1}\bigl|p(x)-|x|\bigr|
=e_K.
\]
These are the required measures and identities.
% lean: Causalean.Mathlib.Analysis.AbsoluteValueMomentPriorDuality.AbsExtremalDecomposition.toMomentMatchedPriors
% lean: CausalSmith.Stat.LdpOptvalueUniformFrontier.bestAbsApproxError_eq_library
% lean: CausalSmith.Stat.LdpOptvalueUniformFrontier.caiLowMomentDuality_proved
\end{enumerate}
\end{proof}

\begin{proof}[Proof of \cref{obj:lem:measurable-kernel-radon-nikodym}]
\begin{enumerate}
\item Use the canonical kernel Radon--Nikodym derivative to define
\[
f:H\times Z\longrightarrow[0,\infty],
\qquad
f(h,z):=\frac{dK}{dL}(h,z).
\]
The kernel Radon--Nikodym construction gives joint measurability of this function for finite kernels with standard-Borel target. % lean: ProbabilityTheory.Kernel.measurable_rnDeriv

\item Fix \(h\in H\) and a measurable set \(E\subseteq Z\). The canonical kernel derivative agrees with the ordinary measure derivative:
\[
f(h,z)
=
\frac{dK(h,\cdot)}{dL(h,\cdot)}(z)
\qquad\text{for \(L(h,\cdot)\)-almost every \(z\)}.
\]
Consequently,
\[
\int_E f(h,z)\,L(h,dz)
=
\int_E
\frac{dK(h,\cdot)}{dL(h,\cdot)}(z)\,L(h,dz)
=
K(h,E).
\]
The first equality uses the almost-everywhere agreement, and the second is the measure Radon--Nikodym identity, applicable because \(K(h,\cdot)\ll L(h,\cdot)\) and both measures are finite. Reversing this equality gives the required identity for every \(h\) and \(E\). % lean: ProbabilityTheory.Kernel.setLIntegral_rnDeriv
\end{enumerate}
\end{proof}
